\documentclass[11pt]{article}

% F1R3FLY.IO branded edition. Brand system: f1r3fly-brand-portal, brand kit v1.3.
% Surface #000000; White, Brand Yellow #F3D630, Brand Sky #3FA9F5;
% League Spartan headings, Source Sans 3 body; gradient accents at hairline weight.

\usepackage[letterpaper,top=1.05in,bottom=1.1in,left=1.05in,right=1.05in,headsep=0.3in,footskip=0.5in]{geometry}
\usepackage{amsmath,amssymb,amsfonts,stmaryrd,mathtools}
\usepackage{amsthm}
\usepackage[no-math]{fontspec}
\usepackage{array}
\usepackage{booktabs}
\usepackage{colortbl}
\newcolumntype{L}[1]{>{\raggedright\arraybackslash}p{#1}}
\usepackage{enumitem}
\usepackage{listings}
\usepackage{xcolor}
\usepackage{pagecolor}
\usepackage{graphicx}
\usepackage{tikz}
\usepackage{fancyhdr}
\usepackage{titlesec}
\usepackage[protrusion=true,expansion=false]{microtype}
\usepackage[hidelinks]{hyperref}

\usetikzlibrary{arrows.meta,positioning,fit,backgrounds,calc}

\emergencystretch=2.5em

% ---------------------------------------------------------------- colour
\definecolor{brandblack}{HTML}{000000}
\definecolor{brandwhite}{HTML}{FFFFFF}
\definecolor{brandyellow}{HTML}{F3D630}
\definecolor{brandsky}{HTML}{3FA9F5}
\definecolor{muted}{HTML}{C5C5C5}
\definecolor{rulegray}{HTML}{959595}
\definecolor{panel}{HTML}{3F3F3F}
\pagecolor{brandblack}
\AtBeginDocument{\color{brandwhite}}
\arrayrulecolor{rulegray}

% ---------------------------------------------------------------- type
\setmainfont{Source Sans 3}[
  Path=fonts/, Extension=.ttf,
  UprightFont=SourceSans3-400, BoldFont=SourceSans3-700,
  ItalicFont=SourceSans3-400, ItalicFeatures={FakeSlant=0.18},
  BoldItalicFont=SourceSans3-700, BoldItalicFeatures={FakeSlant=0.18}]
\setsansfont{Source Sans 3}[
  Path=fonts/, Extension=.ttf,
  UprightFont=SourceSans3-400, BoldFont=SourceSans3-700,
  ItalicFont=SourceSans3-400, ItalicFeatures={FakeSlant=0.18}]
\setmonofont{DejaVu Sans Mono}[Scale=0.84]
\newfontfamily\spartan{League Spartan}[
  Path=fonts/, Extension=.ttf,
  UprightFont=LeagueSpartan-SemiBold, BoldFont=LeagueSpartan-Bold]
\newfontfamily\spartanlight{League Spartan}[
  Path=fonts/, Extension=.ttf, UprightFont=LeagueSpartan-Light]
\newfontfamily\spartanregular{League Spartan}[
  Path=fonts/, Extension=.ttf, UprightFont=LeagueSpartan-Regular]
\newfontfamily\spartanbold{League Spartan}[
  Path=fonts/, Extension=.ttf, UprightFont=LeagueSpartan-Bold]
\newcommand{\eyebrow}[1]{{\spartanbold\addfontfeatures{LetterSpace=12}\MakeUppercase{#1}}}

% ---------------------------------------------------------------- rules
\newcommand{\gradrule}[1][\linewidth]{%
  \noindent\begin{tikzpicture}
    \shade[left color=brandyellow,right color=brandsky] (0,0) rectangle (#1,0.6pt);
  \end{tikzpicture}}

% ---------------------------------------------------------------- headings
\titleformat{\section}
  {\spartan\LARGE\color{brandwhite}}
  {\color{brandyellow}\thesection}{0.6em}{}
  [\vspace{4pt}\gradrule]
\titleformat{\subsection}
  {\spartan\large\color{brandyellow}}
  {\thesubsection}{0.6em}{}
\titleformat{\paragraph}[runin]
  {\spartan\normalsize\color{brandwhite}}{}{0pt}{}
\titlespacing*{\section}{0pt}{3.2ex plus 1ex minus .2ex}{2.4ex plus .2ex}
\titlespacing*{\subsection}{0pt}{2.8ex plus 1ex minus .2ex}{1.4ex plus .2ex}
\titlespacing*{\paragraph}{0pt}{1.6ex plus .5ex}{0.8em}

% ---------------------------------------------------------------- header, footer
\pagestyle{fancy}
\fancyhf{}
\renewcommand{\headrulewidth}{0pt}
\fancyhead[L]{\raisebox{-2pt}{\includegraphics[height=10pt]{logos/f1r3lang-logo-color-on-dark.pdf}}}
\fancyhead[R]{\color{muted}\footnotesize\spartanregular Monadic types and the \texttt{foreach} comprehension}
\fancyfoot[L]{\raisebox{-3pt}{\includegraphics[height=11pt]{logos/f1r3fly-io-bug-icon-color-on-dark.pdf}}\hspace{6pt}%
  {\color{muted}\scriptsize\spartanbold\textcopyright{} F1R3FLY.IO, 2026. All Rights Reserved \textbar{} October 26, 2026}}
\fancyfoot[R]{{\color{muted}\scriptsize\spartanbold\thepage}}
\fancypagestyle{plain}{\fancyhf{}%
  \fancyfoot[L]{\raisebox{-3pt}{\includegraphics[height=11pt]{logos/f1r3fly-io-bug-icon-color-on-dark.pdf}}\hspace{6pt}%
  {\color{muted}\scriptsize\spartanbold\textcopyright{} F1R3FLY.IO, 2026. All Rights Reserved \textbar{} October 26, 2026}}%
  \fancyfoot[R]{{\color{muted}\scriptsize\spartanbold\thepage}}}

% ---------------------------------------------------------------- theorems
\newtheoremstyle{brandplain}{8pt}{8pt}{\itshape}{}{\spartan\color{brandyellow}}{.}{0.5em}{}
\newtheoremstyle{branddef}{8pt}{8pt}{}{}{\spartan\color{brandyellow}}{.}{0.5em}{}
\newtheoremstyle{brandrem}{8pt}{8pt}{}{}{\spartan\color{brandsky}}{.}{0.5em}{}
\theoremstyle{branddef}
\newtheorem{definition}{Definition}[section]
\newtheorem{example}[definition]{Example}
\newtheorem{obligation}[definition]{Obligation}
\theoremstyle{brandplain}
\newtheorem{claim}[definition]{Claim}
\newtheorem{proposition}[definition]{Proposition}
\theoremstyle{brandrem}
\newtheorem{remark}[definition]{Remark}
\newtheorem{aside}[definition]{Aside}
\newtheorem{question}[definition]{Open question}
\renewcommand{\qedsymbol}{\color{brandsky}$\square$}

% ---------------------------------------------------------------- macros
\newcommand{\rhoc}{\textsf{rho}}
\newcommand{\sem}[1]{\llbracket #1 \rrbracket}
\newcommand{\Proc}{\mathsf{Proc}}
\newcommand{\T}{\mathbb{T}}
\newcommand{\M}{\mathsf{M}}
\newcommand{\unit}{\mathsf{unit}}
\newcommand{\fmap}{\mathsf{map}}
\newcommand{\bind}{\mathsf{flatMap}}
\newcommand{\zero}{\mathsf{zero}}
\newcommand{\meas}{\mathsf{measure}}
\newcommand{\alg}{\mathsf{alg}}
\newcommand{\cost}{\mathsf{C}}
\newcommand{\hist}{\mathsf{H}}
\newcommand{\hyp}{\mathsf{Hyp}}
\newcommand{\lax}{\bigcirc}
\newcommand{\Stab}{\mathrm{Stab}}
\newcommand{\red}{\rightsquigarrow}
\newcommand{\nf}{\mathrm{nf}}
\newcommand{\match}{\mathrm{match}}
\newcommand{\chaos}{\mathbf{1}}
\newcommand{\signed}[2]{\{#1\}_{#2}}
\newcommand{\fe}{\texttt{foreach}}
\newcommand{\gate}[2]{\texttt{\{\%}\,#1\,\texttt{\%\}[}#2\texttt{]}}

% The language is f1r3lang: lowercase, always written so (brand kit, section 2).
\newcommand{\flang}{\textsf{f1r3lang}}
\newcommand{\flangbf}{\textsf{\bfseries f1r3lang}}

% ---------------------------------------------------------------- listings
\lstdefinestyle{code}{
  basicstyle=\ttfamily\footnotesize\color{brandwhite},
  keywordstyle=\color{brandsky}\bfseries,
  stringstyle=\color{brandyellow},
  commentstyle=\color{muted}\itshape,
  showstringspaces=false,
  breaklines=true,
  columns=fullflexible,
  frame=leftline,
  rulecolor=\color{brandsky},
  framerule=0.6pt,
  framesep=8pt,
  xleftmargin=12pt,
  aboveskip=10pt,
  belowskip=10pt,
}

\lstdefinelanguage{rholang}{
  morekeywords={new,in,for,foreach,yield,by,match,contract,where,if,then,else,import},
  morecomment=[l]{//},
  morestring=[b]",
}

\lstdefinelanguage{mettail}{
  morekeywords={Module,Theory,theory,import,from,as,Types,Exports,Replacements,
    Terms,Equations,Rewrites,Monads,over,unit,flatMap,map,toBag,ordered,
    commutative,laws,Empty,free,let,in,if,then,subst},
  sensitive=true,
  morecomment=[l]{//},
  morestring=[b]",
}

\lstdefinelanguage{bnf}{
  morekeywords={foreach,yield,where,by},
  sensitive=true,
  morecomment=[l]{--},
  morestring=[b]",
}

\lstset{style=code,language=rholang}

% ---------------------------------------------------------------- abstract
\renewenvironment{abstract}
  {\par\noindent{\color{brandyellow}\eyebrow{Abstract}}\par\vspace{6pt}\small}
  {\par\vspace{6pt}\gradrule\par}

\begin{document}
\begin{titlepage}
\thispagestyle{empty}
\vspace*{0.2in}
\noindent\raisebox{0.55in}{\includegraphics[height=0.42in]{logos/f1r3lang-logo-color-on-dark.pdf}}\hfill\includegraphics[height=1.25in]{logos/f1r3fly-io-vertical-logo-color-on-dark.pdf}

\vspace{1.6in}
\noindent{\color{brandyellow}\eyebrow{Design note \textperiodcentered{} Draft 3}}

\vspace{16pt}
\noindent{\spartanlight\fontsize{38}{45}\selectfont Monadic types and the\\[2pt] \texttt{\fontsize{34}{45}\selectfont foreach} comprehension\\[2pt] in f1r3lang\par}

\vspace{20pt}
\gradrule

\vspace{18pt}
\noindent{\large A synchronous comprehension over declared data monads,\\ with programmable aggregation and explicit cost\par}

\vfill
\noindent{\spartanregular\color{muted}L.G. Meredith\\[2pt] F1R3FLY.io\\[2pt] 26 October 2026}
\end{titlepage}


\begin{abstract}
\noindent
\flang{} (formerly MeTTaIL, and before that rholang~1.4)
has one comprehension, \texttt{for}, and it binds in the \emph{space}: its
sources are channels, its data are shared with concurrent writers, and when
several candidate matchings are available it resolves them by choosing one.
This note adds a second comprehension,
\begin{center}
\texttt{foreach( ptrn <- mndInst where cond ) yield( rslt ) P}
\end{center}
which binds in a \emph{declared data monad}: its source is a term of a monadic
sort, it meets no writer, and it resolves its candidates by aggregating all of
them. Every surviving binding contributes the normal form of its body, the
contributions are combined, and the result is published once, on the channel
\texttt{rslt}, for any process holding that channel to consume.

The design has four parts. In the data definition sublanguage, a
\texttt{Monads} builder declares that a sort is a monad, gives \(\unit\) and
\(\bind\) as rewrites, and supplies a monad morphism \(\mathsf{toBag}\) into
finite bags, whose size is the \emph{measure} that every cost bound in the note
rests on. The monad laws are proof obligations, not equations, so matching
modulo a theory never meets higher-order unification. In the control flow
sublanguage, bodies are evaluated by \emph{lookahead}: each is reduced to normal
form, in isolation, within a terminating fragment picked out by typing --- first
System F, then Abramsky's computational interpretation of linear logic, each
translated into \flang{} with its types translated into spatial-behavioural
formulae. The
comprehension is then a single rewrite whose reduct is one output, which is
the precise sense in which it is synchronous, and which is what entitles it to
stand on the left of \texttt{;}. Aggregation is \emph{programmable}: a
\texttt{by} clause names a theory instance descending from
\texttt{CommutativeMonoid}, checked by origin, and the default is the
\texttt{ParMonoid} of parallel composition. Because parallel composition is the
free commutative monoid, every other aggregator is defined by the universal
property, and accumulation can be tree-shaped without changing the meaning.
In the cost-accounted surface the comprehension carries two gates, an outer
one charged per binding and an inner one funding each body's lookahead. Typing
guarantees termination, not a bound: a comprehension that is not sufficiently
funded does not complete, and its deployer is notified.

A worked example declares the parametric \rhoc{} calculus --- the syntax of
the calculus as a type constructor in its data and its variables --- as a
syntax monad whose bind is semantic substitution, with the group action of
oriented processes as a distributive law and with symmetry reducing the
precharge of a comprehension over a symmetric template.

\fe{} is a conservative extension of the core, so it need not be implemented by
translation into \texttt{for}. The translation is retained as the reference
semantics, and a native implementation is correct when it is \emph{gradedly}
bisimilar to it --- the same behaviour and the same bill.
\end{abstract}

\clearpage
\tableofcontents
\clearpage

%=============================================================================
\section{Two comprehensions}
\label{sec:intro}

The comprehension is the shape data access takes when written down
honestly: name a source, filter it, project the survivors, act. SQL,
LINQ~\cite{LINQ}, XQuery's FLWOR and the list comprehension~\cite{Wadler87}
are the same shape, and Wadler's observation that the comprehension is
monadic~\cite{Wadler92}, in the sense Moggi gave to notions of
computation~\cite{Moggi91}, is what makes the family a family: every one of them
is a notation for \(\bind\) in some monad.

\flang{}'s receipt
\begin{lstlisting}
for( ptrn <- chan where cond ) P
\end{lstlisting}
is that shape with the \texttt{DO} made first-class~\cite{F1R3Dev}. But the
monad it binds in is unusual. The source is a channel; the datum is produced by
a concurrently running writer; the join \texttt{\&} is a rendezvous; the
communication rule substitutes the datum into the continuation and thereby
\emph{witnesses} the meeting. Call this the \emph{space} monad. Its bind is
asynchronous, it is shared, and it is nondeterministic in the strong sense that
competing readers race.

Developers also need the other kind of comprehension: the one that LINQ, and
Scala's \texttt{for} with \texttt{yield}, provide over a collection already in hand. A
term received on a channel may itself be a list, a bag, a tree of findings, a
distribution. Walking it with \texttt{for} is possible only by first spilling
it onto a channel, which is both expensive and wrong: it exposes a private
traversal to every other reader of that channel, and it turns a deterministic
computation over a value into a race.

\fe{} is the missing comprehension. It binds in a \emph{data monad} declared in
the data definition sublanguage (DDL), evaluates a body for every surviving
binding, combines the results, and publishes the combination once:
\begin{lstlisting}
foreach( i <- { 1, 3, 5, 7, 9 } where i > 4 ) yield( rslt ) i
  ~>  rslt!( 5 | 7 | 9 )
\end{lstlisting}
The result is on a channel because a result is something other processes
consume, and in \flang{} that is what channels are for. A consumer may take it,
persist on it, or peek at it with the three arrows of \texttt{for}.

\paragraph{Choose versus sum.} The cleanest way to see the difference is through
the word the developer document uses for the thread running through
\flang{}'s semantics: \emph{candidate}~\cite{F1R3Dev}. A candidate is one way
of assigning data to the patterns of a comprehension.
\begin{itemize}[nosep]
\item In \texttt{for}, candidates are resolved by \emph{choice}. One fires; the
others do not.
\item In \fe{}, candidates are resolved by \emph{aggregation}. Every candidate
contributes --- a rejected one contributes the aggregator's unit --- and the
contributions are combined.
\end{itemize}
Figure~\ref{fig:choosesum} draws the distinction. Section~\ref{sec:graded}
shows that it is exactly the distinction the graded \texttt{where} clause needs.

\begin{figure}[t]
\centering
\resizebox{\linewidth}{!}{%
\begin{tikzpicture}[color=white,
  node distance=6mm and 10mm,
  box/.style={draw, rounded corners=2pt, minimum width=14mm, minimum height=6mm,
    font=\footnotesize\ttfamily},
  lab/.style={font=\footnotesize},
  >={Stealth[length=2mm]}]
  % left: for
  \node[box] (r) {for(p<-x)};
  \node[box, above right=of r] (c1) {c1};
  \node[box, right=of r] (c2) {c2};
  \node[box, below right=of r] (c3) {c3};
  \draw[->, thick] (r) -- (c2);
  \draw[->, dashed, rulegray] (r) -- (c1);
  \draw[->, dashed, rulegray] (r) -- (c3);
  \node[box, right=of c2] (k1) {P\{c2\}};
  \draw[->, thick] (c2) -- (k1);
  \node[lab, below=12mm of r] {\texttt{for}: one candidate is chosen};
  % right: foreach
  \begin{scope}[xshift=78mm]
  \node[box] (m) {foreach(p<-m)};
  \node[box, above right=of m] (d1) {c1};
  \node[box, right=of m] (d2) {c2};
  \node[box, below right=of m] (d3) {c3};
  \foreach \d in {d1,d2,d3} \draw[->, thick] (m) -- (\d);
  \node[box, right=of d2, minimum width=34mm] (k2) {rslt!(nf(P\{c1\}) | ... )};
  \foreach \d in {d1,d2,d3} \draw[->, thick] (\d) -- (k2);
  \node[lab, below=12mm of m] {\fe{}: every candidate contributes};
  \end{scope}
\end{tikzpicture}}
\caption{The two resolution disciplines for candidates. Under the default
aggregator the published result is a parallel composition, because
\texttt{PPar} is a bag (\S\ref{sec:free}).}
\label{fig:choosesum}
\end{figure}

\paragraph{Contributions.}
\begin{enumerate}[nosep]
\item A \texttt{Monads} builder for the DDL, with the laws as obligations and a
monad morphism into finite bags that both drives the comprehension and measures
it (\S\ref{sec:ddl}).
\item The \fe{} comprehension: its syntax, its meaning as a single rewrite,
evaluation of bodies by lookahead in a terminating fragment given by System F
and then by Abramsky's proof expressions, and a reference translation into the
core (\S\ref{sec:cfl}).
\item Programmable aggregation: aggregators as theory instances descending
from the universal-algebra theory of commutative monoids, defined by the universal property of parallel
composition, accumulated as a tree, and extended in three further stages
(\S\ref{sec:agg}).
\item Sequencing with \texttt{;}, on the strength of termination
(\S\ref{sec:seq}).
\item An explicit cost-accounted surface with two gates, per-binding charging,
an atomic outer gate, and precharge estimates, with non-completion and
notification as the policy for underfunded runs
(\S\ref{sec:cost}).
\item An implementation stance --- native evaluation, the translation retained
as specification, graded bisimulation as the correctness criterion --- and a
map of what changes in the node (\S\ref{sec:impl}).
\item An account of the interaction with the rest of \flang{}
(\S\ref{sec:interactions}), a worked example on the parametric \rhoc{} calculus
(\S\ref{sec:rho}), and a first tranche (\S\ref{sec:tranche}).
\end{enumerate}

%=============================================================================
\section{Background}
\label{sec:background}

This section fixes the parts of \flang{} the feature touches. Fuller
treatments are in the developer document~\cite{F1R3Dev} and the GSLT
reference~\cite{GSLTintro}. \flang{} is built on the reflective higher-order
\rhoc{} calculus~\cite{MR05}, presented in its own DDL as one theory among
others.

\subsection{The data definition sublanguage}

A language definition is a \texttt{Theory}. Its body is an ordered chain of
postfix builders --- \texttt{Types}, \texttt{Exports}, \texttt{Replacements},
\texttt{Terms}, \texttt{Equations}, \texttt{Rewrites} --- applied to a starting
theory, and theories combine by join, meet and difference over shared ancestors
\cite{modulesyntax}. Every element introduced by a builder carries an origin
token, and two elements are identified exactly when their origins agree; this is
what makes a join a pushout rather than a union, and what produces the
\texttt{[join-collision]} and \texttt{[unshared-name]} diagnostics. A
replacement renames an element and keeps its origin.

Term constructors are written in judgement form,
\begin{lstlisting}[language=mettail]
PInput . n:Name, ^x.p:[Name -> Proc] |- n "?" x "." "{" p "}" : Proc;
PPar   . ps:HashBag(Proc) |- "{" ps.*sep("|") "}" : Proc;
\end{lstlisting}
with binders \verb|^x.p| as first-class arguments, and three built-in
collection sorts \texttt{HashBag}, \texttt{Set} and \texttt{List} over a
category. There are no user-declared sort constructors: in the reference
elaborator a sort is either a category or a built-in collection of one. That
fact governs the first tranche of this design.

The reference corpus includes a universal-algebra module whose tower of
theories, from \texttt{EmptySet} and \texttt{Monoid} through
\texttt{CommutativeMonoid} and \texttt{Group} to \texttt{Rig} and
\texttt{Ring}, is built
by parameters and replacements, and the \rhoc{} calculus's own parallel
composition is defined as an instance of it:
\begin{lstlisting}[language=mettail]
Theory ParMonoid(cm: u.CommutativeMonoid) {
  cm
    Exports      { Elem => Proc; }
    Replacements { Zero => PZero . |- "0" : Proc;
                   Plus => PPar  . ps:HashBag(Proc) |- "{" ps.*sep("|") "}" : Proc; }
}
\end{lstlisting}
That definition is the hinge of \S\ref{sec:agg}.

\subsection{The control flow sublanguage}

The comprehension has the form
\begin{lstlisting}
for( ptrn_11 <- chan_11 & ... & ptrn_1n <- chan_1n  where cond_1 ;
     ...                                                         ;
     ptrn_k1 <- chan_k1 & ... & ptrn_km <- chan_km  where cond_k ) P
\end{lstlisting}
with \texttt{\&} for simultaneity within a group, \texttt{;} for sequence
between groups, and a \texttt{where} per group that may mention any variable
bound in that group or an earlier one. Three arrows distinguish a consuming
read \texttt{<-}, a persistent read \texttt{<=} and a peek \texttt{<<-}. A
failed guard leaves the datum in the space: a failed transaction is not rolled
back, it never began~\cite{F1R3Dev}.

Rholang~1.1 also provides a synchronous send, \verb|x!?(Q); K|, which the
node's normaliser desugars as
\begin{lstlisting}
new r in { x!(*r, Q) | for( _ <- r ){ K } }
\end{lstlisting}
so that \texttt{K} runs once the receiver on \texttt{x} acknowledges on the
fresh name \texttt{r}~\cite{f1r3nodeOTE}. Chains of the form
\verb|x1!?(Q1); ...; xn!?(Qn); K| are nested instances. The continuation runs
only if the receiver chooses to reply.

\subsection{The cost-accounted surface}

Cost accounting is written in the language rather than bolted onto an
interpreter~\cite{CostRho}. A signed term \(\signed{P}{s}\), concretely
\verb|{% P %}[s]|, runs \texttt{P} once a token for signature \texttt{s} is
available; a budget is a stack \verb|s :: s :: ... :: ()| whose depth is the
number of operations it funds. Signatures are ground names, sections
\verb|# P|, and compounds. The surface already admits per-clause signed binds,
\begin{lstlisting}
{% for( {% y1 <- x1 %}[s1] & {% y2 <- x2 %}[s2] ) P %}[ s0 ]
\end{lstlisting}
in which each rendezvous surface is funded independently, and in the native
four-sort grammar the continuation and payload slots of \texttt{for} and send
have the signed-term sort \(\T\). That sorting --- \emph{wrapping by
construction} --- is what makes subject reduction coincide with the absence of
unmetered execution~\cite{CostMonad}. Funded runs are bounded: strong
normalisation holds conditionally on funding~\cite{f1r3nodeOTE}. That is a
statement about budgets, not types, and \S\ref{sec:lookahead} keeps the two
apart.

\subsection{Monads one level up}

The GSLT reference treats three monads: the cost monad \(\cost\), the history
monad \(\hist\) and the Hypercube construction \(\hyp\)~\cite{GSLTintro,
CostMonad, History, Hypercube}. All three are endofunctors on the category of
theories: they take a theory and return a metered, logged or typed theory.

The monads of this note live one level down. A declaration
``\(\M\) is a monad'' is a statement \emph{inside} a theory, about one of its
sorts. The distinction matters in two places. In prose, the phrase ``the
monad'' becomes ambiguous; this note says \emph{data monad} for the new kind and
\emph{theory-level monad} for \(\cost\), \(\hist\) and \(\hyp\). And in
semantics, the theory-level monads act on theories that contain data monads,
which raises the question of whether the declaration survives them.

\begin{proposition}[Laws modulo erasure]
\label{prop:modulo-erasure}
Let theory \(T\) declare a data monad \(\M\) whose laws hold in \(T\). Then
\(\cost(T)\) and \(\hist(T)\) carry the declaration, and the laws hold in each
up to the corresponding erasure back to \(T\), but not in general on the nose.
\end{proposition}

\begin{proof}[Sketch]
The two sides of the left unit law,
\(\bind(\unit(a), k)\) and \(k\,a\), differ by one application of the
\(\bind\) rewrite. In \(\cost(T)\) that application is gated and consumes a
token; in \(\hist(T)\) it is logged as an event. The two sides therefore differ
in cost and in history, and agree after the erasures \(\cost(T)\to T\) and
\(\hist(T)\to T\) studied in~\cite{TwoErasures}. The right unit and
associativity laws are similar.
\end{proof}

This is the expected behaviour, not a defect, and it anticipates the
correctness criterion of \S\ref{sec:impl}: once cost is explicit, ``the same
computation'' means the same behaviour \emph{and} the same bill.

%=============================================================================
\section{Declaring a data monad}
\label{sec:ddl}

\subsection{The prerequisite: a sort constructor}

A monad is a sort constructor \(\M(-)\) with structure. \(\fmap\) alone needs
two element sorts, \((A \to B) \to \M(A) \to \M(B)\). The DDL has no
user-declared sort constructors, so the design proceeds in two steps.

\paragraph{Step 1: the built-in collections.} \texttt{List}, \texttt{Set} and
\texttt{HashBag} are already sort constructors, and they are the canonical
monads: the list monad, the finite powerset monad, and the finite multiset
monad. Declaring them monadic needs no new sort machinery. They are declared once,
in a prelude supplied by the toolchain, and are available to every theory.

\paragraph{Step 2: sort parameters.} The general feature extends
\texttt{Types} with parameterised categories,
\begin{lstlisting}[language=mettail]
Types { Tree(A); }
\end{lstlisting}
where \texttt{A} ranges over categories. This is the real feature, and it is
independently useful: the Hypercube construction currently generates type
systems over unparameterised sorts~\cite{Hypercube}. Its motivating instance is
the calculus itself, whose syntax is already treated in the working notes as a
type constructor in its data and its variables; \S\ref{sec:rho} works it
through.

\paragraph{A rejected alternative.} A theory parameter can stand in for a type
parameter --- \texttt{Theory List(e: EmptySet)} is a list of whatever
\texttt{e} exports. But then \(\M\) is a map from theories to theories, the
feature climbs back to the level of \(\cost\) and \(\hist\), and the sharing
discipline objects: \texttt{List(List(s))} reintroduces the list constructors
along its own chain and draws \texttt{[repeat-label]}. The diagnostic is
correct, and the alternative is rejected.

\subsection{The \texttt{Monads} builder}

A new postfix builder declares one or more data monads. Its entries name the
monad, the sort it is a monad on, the interpretation of the operations, and
the structure it offers to the comprehension. The prelude declarations of the
list and multiset monads are written with a sort variable \texttt{A} that
Step~2 makes available to users and that the prelude may use from the start.

\begin{lstlisting}[language=mettail]
Monads {
  Seq over List {
    unit    : (a)                    ~> [a];
    flatMap : ([], ^x.k)             ~> [];
    flatMap : ([a, ...r], ^x.k)      ~> (Concat (subst ^x.k a) (flatMap r ^x.k));
    toBag   : ([a, ...r])            ~> {a, ...(toBag r)};
    ordered;
    laws    : LeftUnit, RightUnit, Assoc, ToBagMorphism;
  }
  Bag over HashBag {
    unit    : (a)                    ~> {a};
    flatMap : ({}, ^x.k)             ~> {};
    flatMap : ({a, ...r}, ^x.k)      ~> (Union (subst ^x.k a) (flatMap r ^x.k));
    toBag   : identity;
    commutative;
    laws    : LeftUnit, RightUnit, Assoc, Comm, ToBagMorphism;
  }
}
\end{lstlisting}

The surface is illustrative; its content is not. Five decisions are embedded
in it.

\paragraph{The interpretation is a set of rewrites.} The DDL defines
operations only through rules, so \(\unit\) and \(\bind\) are given by cases
over the sort's constructors, with the continuation as a first-class
abstraction \verb|^x.k| of sort \verb|[A -> M(B)]|. These are ordinary
rewrites: they are oriented, they fire, and they ``wiggle'' like any other
rule of the theory. \(\fmap\) may be given explicitly for efficiency; when it
is, \(\fmap\,f = \bind(-, \unit\circ f)\) becomes one more obligation.

\paragraph{The laws are obligations, not equations.} Matching in \flang{} is
matching modulo the equations of the theory that types the channel. Placing
\(\bind(m, \lambda x.\unit(x)) = m\) in an \texttt{Equations} block would put
the matcher into higher-order unification, which is undecidable in general.
The \texttt{laws} entry therefore \emph{names} the laws; the elaborator records
them as obligations, discharged by proof (as the cost monad's laws are, in
Rocq and Lean) or by bounded exhaustive checking (as in the Sage harness), and
never hands them to the matcher.

\begin{obligation}[Monad laws]
For every declared data monad, the left unit, right unit and associativity
laws hold up to the equations of the theory.
\end{obligation}

\begin{obligation}[Confluence with the host theory]
\label{obl:confluence}
The \(\bind\) and \(\mathsf{toBag}\) rewrites have no critical pair with the
other rewrites of the theory, modulo its equations, that fails to join.
\end{obligation}

\paragraph{\(\mathsf{toBag}\) is what the comprehension uses.} A comprehension
visits the elements of its source with their multiplicities. That is a natural
transformation \(\mathsf{toBag} : \M \Rightarrow \mathsf{HashBag}\), the
support with multiplicity, and it is the one piece of structure every monad
used with \fe{} must supply. For dependent groups (\S\ref{sec:groups}) to agree
with the monad's own sequencing it must be a monad morphism:
\[
\mathsf{toBag}(\unit(a)) = \{a\},
\qquad
\mathsf{toBag}(\bind(m, k)) = \textstyle\biguplus_{x \in \mathsf{toBag}(m)} \mathsf{toBag}(k\,x).
\]
It is one for \texttt{Seq} (forget order) and \texttt{Bag} (identity). It is
not one for the finite powerset, whose flatten removes duplicates; for
\texttt{Set} the comprehension therefore visits \emph{binding instances}, and
deduplication of results is a choice of aggregator, not of source
(\S\ref{sec:stages}).

\begin{obligation}[Bag morphism]
\label{obl:tobag}
\(\mathsf{toBag}\) is a monad morphism \(\M\Rightarrow\mathsf{HashBag}\).
\end{obligation}

The \emph{measure} of an instance is the size of its bag,
\[
\meas(m) \;=\; |\mathsf{toBag}(m)|,
\]
and it is the quantity on which the completion of the reference translation
(\S\ref{sec:refsem}) and every cost bound of \S\ref{sec:cost} rest. Defining it
this way, rather than declaring it separately, makes the earlier ``exact
measure'' obligation a consequence of Obligation~\ref{obl:tobag}: \(\bind(m,
k)\) invokes \(k\) once per element of \(\mathsf{toBag}(m)\).

\paragraph{Capabilities are claims, and each one licenses a feature.} Beyond
\(\mathsf{toBag}\), two optional entries matter, and each is checked where it is
used.

\begin{center}
\small
\begin{tabular}{@{}L{0.17\textwidth}L{0.33\textwidth}L{0.40\textwidth}@{}}
\toprule
\textbf{Entry} & \textbf{Meaning} & \textbf{Licenses} \\
\midrule
\texttt{toBag} (required) & elements with multiplicity; a monad morphism &
\fe{} over the monad at all; its measure \\[3pt]
\texttt{ordered} & a canonical enumeration order on \(\mathsf{toBag}(m)\) &
non-commutative (ordered) aggregators (\S\ref{sec:stages}) \\[3pt]
\texttt{commutative} & order of independent binds is immaterial~\cite{Kock70} &
the join \texttt{\&} under an ordered aggregator (\S\ref{sec:groups}) \\
\bottomrule
\end{tabular}
\end{center}

\noindent
Two capabilities the earlier draft required have gone. A \(\zero\) is no longer
needed for filtering, because a rejected candidate contributes the
\emph{aggregator's} unit (\S\ref{sec:where}); and an algebra on processes is no
longer needed, because the result is data published on a channel, and
aggregation is the job of the aggregator (\S\ref{sec:agg}).

\paragraph{One sort may carry several monads.} The monads are named
(\texttt{Seq}, \texttt{Bag}) because a sort can legitimately carry more than
one monad structure, just as a rig carries two monoids on one carrier. A
comprehension may name the monad it binds in, \texttt{foreach[Seq](...)};
the qualifier may be omitted when exactly one monad is declared on the
source's sort.

\subsection{Monads and the algebra of theories}
\label{sec:modules}

A \texttt{Monads} entry is an element like any other, so it receives an origin
token at introduction and is carried through inheritance.

\begin{itemize}[nosep]
\item \textbf{Join.} Two theories that inherit \texttt{Seq} from one shared
instance join cleanly, and the comprehension sees one monad. Two theories that
independently declare a monad under the same name on the same sort draw
\texttt{[join-collision]}: they are two structures with one spelling.
\item \textbf{Difference.} Subtracting a constructor that a monad's rules
mention leaves the declaration dangling. That is a new diagnostic,
\texttt{[dangling-monad]}, and it is the analogue of the existing
\texttt{[forward-reference]}.
\item \textbf{Ordering.} A \texttt{Monads} block must follow the \texttt{Terms}
blocks whose labels its rules mention, by the existing chain discipline.
\item \textbf{Replacements.} Renaming a constructor preserves its identity and
therefore preserves the monad. Replacing a constructor by one of different arity
reproduces the open arity-drift problem already recorded for \texttt{PPar}:
the inherited \(\bind\) rules are written over a shape that no longer exists.
Whatever resolution is adopted there --- arity-preserving replacements, or a way
to declare inherited rules structural --- must be adopted here as well.
\item \textbf{Parameters.} Once Step~2 lands, monad transformers are theories
that take a monad as a parameter. This is the fourth rung of the DDL doing what
it was built for.
\end{itemize}

%=============================================================================
\section{The \texttt{foreach} comprehension}
\label{sec:cfl}

\subsection{Syntax}

\begin{lstlisting}[language=bnf]
Proc     ::= ... | Foreach | Foreach ";" Proc
Foreach  ::= "foreach" Qual? "(" Groups ")" "yield" "(" Name ")" By? Proc
Qual     ::= "[" MonadName "]"
Groups   ::= Group ( ";" Group )*
Group    ::= Bind ( "&" Bind )* ( "where" Cond )?
Bind     ::= Pattern "<-" Term          -- Term of a monadic sort
By       ::= "by" TheoryInstance        -- default: ParMonoid
\end{lstlisting}

The surface coincides with \texttt{for} wherever the meanings coincide and
diverges visibly where they do not. The channel named by \texttt{yield} is free
in the comprehension and is where the result appears. The body \texttt{P} is in
the scope of the pattern variables. Its cost-accounted form is the subject of
\S\ref{sec:cost}; until then, gates are elided.

\subsection{Meaning in one rule}
\label{sec:rule}

For a single group,
\begin{align*}
&\texttt{foreach( }p \leftarrow m\texttt{ where }c\texttt{ ) yield( }r\texttt{ ) by }A\ P\\
&\qquad\red\;
r!\Big(\,\textstyle\bigoplus_{A}\ \big\{\, \mathsf{contrib}(x) \;:\; x \in \mathsf{toBag}(m) \,\big\}\Big)
\end{align*}
where \(\bigoplus_A\) is the operation of the aggregator \(A\) (empty sum: its
unit \(\varepsilon_A\)), and
\[
\mathsf{contrib}(x) \;=\;
\begin{cases}
\nf(P\sigma) & \text{if } \sigma = \match(p, x) \text{ exists and } \nf(c\sigma) = \texttt{true},\\
\varepsilon_A & \text{otherwise.}
\end{cases}
\]
\(\nf\) is the normal form computed by lookahead (\S\ref{sec:lookahead}). Under
the default aggregator \(A = \texttt{ParMonoid}\), \(\bigoplus\) is \texttt{|}
and \(\varepsilon\) is \texttt{0}:
\begin{lstlisting}
foreach( i <- { 1, 3, 5, 7, 9 } where i > 4 ) yield( rslt ) i         ~>  rslt!( 5 | 7 | 9 )
foreach( i <- { 1, 3, 5, 7, 9 } where i > 4 ) yield( rslt ) i * i     ~>  rslt!( 25 | 49 | 81 )
foreach( i <- { 1, 3 }          where i > 4 ) yield( rslt ) i         ~>  rslt!( 0 )
\end{lstlisting}
Three readings of the rule are worth stating at once.
\begin{itemize}[nosep]
\item \emph{Bodies run.} Each surviving body is reduced to normal form, so its
internal computation happens. What it cannot do during the comprehension is
meet a process outside itself; a send or receipt on an outside channel remains,
unreduced, in its normal form and so in the published result.
\item \emph{The result is data.} The publication is a send, so its payload is
quoted, as every send's is. A consumer that wants the contributions to run says
so: \verb|for( @ps <- rslt ){ ps }|.
\item \emph{One step.} From outside, the comprehension is a single reduction
whose reduct is a single output. This is the strongest reading of
``synchronous'' available, and \S\ref{sec:sync} shows that the reference
translation realises it.
\end{itemize}
Dependent groups compose the rule through \(\bind\) (\S\ref{sec:groups}).

\subsection{Static discipline}
\label{sec:static}

\paragraph{Sources are terms of a monadic sort.} In \texttt{for}, a source is a
channel. In \fe{}, a source is a term whose sort carries a declared monad. A
channel may not appear as a \fe{} source; a monadic term may not appear as a
\texttt{for} source except by being quoted into a name, in which case it is a
channel and nothing more. The two comprehensions nest freely in each other's
bodies. This requires something of the CFL it has not needed before: the sort
of the source expression must be known, by inference from the pattern or the
channel type that bound it, or from an explicit qualifier.

\paragraph{One arrow.} Only \texttt{<-} is meaningful. Persistence \texttt{<=}
has no meaning over a value, and a peek \texttt{<<-} is moot because a term
source is never consumed. Both are refused rather than reinterpreted.

\paragraph{One matcher per element.} Matching in \flang{} is modulo the
equations of the theory, and a pattern can have several matchers against one
element (a remainder pattern over a bag, say). The rule uses the matcher the
\texttt{match} construct selects, so each element contributes exactly once. A
pattern with several matchers modulo the equations makes that selection
nondeterministic and places the comprehension outside the deterministic
fragment of \S\ref{sec:lookahead}.

\paragraph{Typing.} Write \(\mathcal{T}\) for the terminating fragment of
\S\ref{sec:lookahead} and \(|A|\) for the carrier sort of the aggregator.
\(\Gamma \vdash_{\mathcal{T}} P : S\) means that \(P\) is \(\sem{M}\) for a
source term \(M\) well typed at a type whose ground translation is \(S\); the
element sort of the monad, bound into \(\Delta\), must likewise be the
translation of a source type.
\[
\frac{
\Gamma \vdash m : \M(B) \quad
p : B \leadsto \Delta \quad
\Gamma,\Delta \vdash_{\mathcal{T}} c : \textsf{Bool} \quad
\Gamma,\Delta \vdash_{\mathcal{T}} P : |A| \quad
\Gamma \vdash r : \textsf{chan}(|A|)
}{
\Gamma \vdash \texttt{foreach[}\M\texttt{]( } p \leftarrow m \texttt{ where } c \texttt{ ) yield( } r \texttt{ ) by } A\ P : \Proc
}
\]
with \(\M\) supplying \(\mathsf{toBag}\) and \(A\) a commutative-monoid
instance (\S\ref{sec:agg}). Under the default aggregator \(|A| = \Proc\), so
every body in \(\mathcal{T}\) is well typed. Since a channel's type is a
theory, \(\textsf{chan}(|A|)\) means a channel typed by a theory that exports
the aggregator's carrier.

\subsection{Groups: \texttt{;} and \texttt{\&}}
\label{sec:groups}

Between groups, \texttt{;} is dependent bind: in
\begin{lstlisting}
foreach( x <- m ; y <- n(x) ) yield( r ) P
\end{lstlisting}
the source of the second group may mention \texttt{x}, and the candidates are
the elements of
\[
\mathsf{toBag}\big(\bind(m,\ \lambda x.\,\fmap(n(x),\ \lambda y.\,(x,y)))\big)
\;=\;
\textstyle\biguplus_{x \in \mathsf{toBag}(m)}\ \{\, (x,y) : y \in \mathsf{toBag}(n(x)) \,\},
\]
the equality being Obligation~\ref{obl:tobag}. The per-group scoping of
\texttt{where} carries over unchanged.

Within a group, \texttt{\&} means \emph{independent} binds: neither source
mentions the other's variables. Under a commutative aggregator --- every
aggregator of Stage~1 --- the order in which independent binds are taken cannot
show in the result, so \texttt{\&} is always permitted. Under an ordered
aggregator (\S\ref{sec:stages}) the order would show, and \texttt{\&} is then
permitted only on monads that claim \texttt{commutative}. The restriction has
moved from the monad to the aggregator.

\subsection{\texttt{where} in \texttt{foreach}}
\label{sec:where}

A rejected candidate contributes the aggregator's unit. Filtering therefore
needs nothing of the monad: no zero, no MonadPlus structure. The diagnostic the
earlier draft proposed for monads without a zero is withdrawn.

The same keyword now has two operational readings, and the developer document
should state both:
\begin{itemize}[nosep]
\item in \texttt{for}, a failed guard \emph{leaves the datum in the space};
\item in \fe{}, a failed guard \emph{contributes the unit}.
\end{itemize}
What reconciles them is that in both cases the rejected candidate is a
non-event: nothing is consumed, nothing is undone, and nothing downstream can
distinguish it from its absence.

\subsection{Lookahead and the terminating fragment}
\label{sec:lookahead}

The rule needs \(\nf(P\sigma)\) and \(\nf(c\sigma)\) to exist and to be unique.
Bodies and guards are therefore evaluated by \emph{lookahead}: each is reduced
to normal form by itself, in isolation from the rest of the configuration, and
the normal form is what is contributed. Existence and uniqueness come from a
type system.

\paragraph{Termination comes from a typed source calculus.} Several type
systems guarantee termination, and they are used the same way. Take a source
calculus \(L\) whose well-typed terms are strongly normalising, together with
two translations: of terms, \(M \mapsto \sem{M}\), into \flang{}; and of types,
\(\tau \mapsto \varphi_\tau\), into spatial-behavioural formulae. The
terminating fragment is the image of the well-typed terms,
\[
\mathcal{T}_L \;=\; \{\, \sem{M} \;:\; \vdash_L M : \tau \,\}.
\]
The division of labour is the point of the construction. The formula
\(\varphi_\tau\) is the \emph{specification} of the fragment: it implies
termination, which in a logic with least fixed points is \(\mu X.\,\Box X\).
The type derivation in \(L\) is a \emph{decidable certificate} for it. Checking
\(\mu X.\,\Box X\) directly against an arbitrary process is undecidable over a
Turing-complete theory; checking a derivation is not.

\paragraph{First fragment: System F.} Well-typed terms of System
F~\cite{GirardLafontTaylor89, Reynolds74} are strongly normalising, by Girard's
reducibility candidates, and the calculus is Church--Rosser. System F is
declared as a \texttt{Theory} in the DDL --- types as a sort, terms with type
abstraction and application, and the two \(\beta\)-rules as rewrites --- and in
the first tranche the guards and bodies of a \fe{} are terms of that theory.
They are type-checked at compile time and compiled by \(\sem{-}\). Annotations
are explicit, since type inference for System F is
undecidable~\cite{Wells99}. Polymorphism is what makes the fragment worth
having: one body can be generic in the element sort of the monad, and an
aggregator's operation can be written once for every carrier it serves.
\begin{lstlisting}
foreach( i <- { 1, 3, 5, 7, 9 } where F( i > 4 ) ) yield( rslt ) by a.IntSum()
  F( (/\a. \f:a->a. \x:a. f (f x)) [Int] (\n:Int. n * n) i )
  ~>  rslt!( 9587 )                       // 5^4 + 7^4 + 9^4
\end{lstlisting}
The body applies a polymorphic \emph{twice} at \texttt{Int}; its normal form at
each survivor is \(i^4\), and the aggregator sums them.

\paragraph{Second fragment: Abramsky's proof expressions.} Well-typed proof
expressions of Abramsky's computational interpretation of linear
logic~\cite{Abramsky94} are terminating as well, and its descendants in the
session-typed tradition are confluent and deadlock-free~\cite{CairesPfenning10,
Wadler12}. Its types land in a different logic from System F's: the linear
connectives are the connectives of the logic generated with quantales as the
container of witnesses~\cite{F1R3Dev}. What it adds is \emph{internal
communication}: a body may consist of processes that talk to one another on
linear channels, and still terminate with a unique normal form --- concurrency
inside a single contribution, not only between contributions. It is also the
natural home of the linear bodies the quantum reading requires
(\S\ref{sec:interactions}). The rho side of the translation is partly worked
out already, since the repair of the \(\pi\)-into-\rhoc{} encoding realises half
of Abramsky's construction~\cite{Namespaces}.

\begin{obligation}[Faithful translation]
\label{obl:fragment}
For each source calculus \(L\) admitted as a fragment:
\begin{enumerate}[nosep,label=(\roman*)]
\item \emph{Divergence is reflected.} Every reduction of \(\sem{M}\) corresponds
to zero or more reductions of \(M\), with finitely many administrative steps
between them, so termination transfers. Lookahead helps: the correspondence is
needed only for a closed term reduced by itself, never under an arbitrary
context.
\item \emph{Confluence transfers.} The translation fixes one reduction strategy,
so the image of a confluent calculus has normal forms unique up to the
equations of the theory.
\item \emph{The formula map is sound.} \(\vdash_L M : \tau\) implies
\(\sem{M} \models \varphi_\tau\), and \(\varphi_\tau\) implies
\(\mu X.\,\Box X\).
\item \emph{Ground values read back.} At base types the translation is the
identity on \flang{}'s ground values, so a body typed \texttt{Int} normalises to
an integer the aggregator can combine.
\end{enumerate}
\end{obligation}

\noindent
Part (iii) asks something of the generated logics: least fixed points, and,
for System F, quantification over formulae to interpret \(\forall\), as
reducibility candidates do. That is more than the spatial-behavioural logics of
the working notes currently provide, and it is needed only for the
\emph{statement} of the fragment; the first tranche rests on (i), (ii), (iv)
and the source typing (Question~\ref{q:logic}).

\paragraph{Typing guarantees termination, not a bound.} System F normalisation
terminates, but its length is not bounded by any elementary function, and the
same is true of linear logic with exponentials. Fragments with a precomputed
cost bound exist --- the multiplicative fragment, light and soft linear logics
--- and were set aside as too constrained to be interesting. So termination is
the type system's responsibility and completion within budget is the
deployer's. Lookahead is metered like everything else (\S\ref{sec:cost}); a
comprehension that is not sufficiently funded does not complete, publishes
nothing, and its deployer is notified. That is the whole policy.

Lookahead is what makes bodies \emph{run} without making the comprehension
wait on the world: a body cannot block on an outside writer, because during
lookahead there is no outside.

\subsection{Reference semantics}
\label{sec:refsem}

The rule is specified by a translation into the core, which is the standard
against which any implementation is judged (\S\ref{sec:impl}). The source is
spilled onto a private channel via \(\mathsf{toBag}\), every element is
visited by a persistent receipt and contributes either the lookahead of its
body or the unit, and contributions are folded into a private accumulator that
publishes once, when a count reaches the measure:
\begin{lstlisting}
[[ foreach( ptrn <- m where cond ) yield( rslt ) P ]] =
  new src, acc, ans in {
    [[ m ]]_src                                    // src!(x1) | ... | src!(xn),  n = measure(m)
  | acc!( 0, Nil )
  | if ( n == 0 ) { rslt!( Nil ) }                 // empty source: publish the unit
  | for( @x <= src ){
      match x {
        ptrn where cond => ans!( [[ P ]] )[*]       // [*]: lookahead to normal form
        _               => ans!( Nil )
      }
    }
  | for( @(k, a) <= acc & @e <= ans ){
      if ( k + 1 == n ) { rslt!( a | e ) } else { acc!( k + 1, a | e ) }
    }
  }
\end{lstlisting}
Three features of the translation carry its correctness.
\begin{itemize}[nosep]
\item \textbf{Initialisation.} The accumulator starts at the unit,
\texttt{acc!( 0, Nil )}; without it no join fires.
\item \textbf{Completion.} The accumulator is private, and \texttt{rslt} is
written exactly once, when the count reaches \(n\). A process holding
\texttt{rslt} can never observe or steal a partial sum, and the remaining folds
can never be stranded.
\item \textbf{Every element counts.} The pattern sits in a \texttt{match} with a
catch-all rather than on the receipt, so an element that fails the pattern or
the guard still contributes, and the count reaches \(n =\meas(m)\) exactly.
\end{itemize}
The measure therefore enters the \emph{semantics}, not only the cost: it is the
completion condition. For dependent groups the total is not known at entry, and
the count becomes a deficit in \(\mathbb{Z}\): it starts at \(\meas(m)\), each
first-group binding \(x\) adds \(\meas(n(x)) - 1\) as it is expanded, each leaf
contribution subtracts one, and publication happens at zero. Integers under
addition are a commutative monoid, so the deficit is one more aggregator
running alongside the first (\S\ref{sec:tree}).

The translation exhibits what the monad must supply: \(\mathsf{toBag}\), and
nothing else, for a single group; \(\bind\), through Obligation~\ref{obl:tobag},
for dependent ones.

\subsection{Synchrony}
\label{sec:sync}

\begin{proposition}[Synchrony]
\label{prop:sync}
Suppose the guards and bodies of a \fe{} are typed in \(\mathcal{T}\). Then its
reference translation is weakly bisimilar to the single output
\(r!(\bigoplus_A\{\mathsf{contrib}(x)\})\) of \S\ref{sec:rule}.
\end{proposition}

\begin{proof}[Sketch]
Every communication of the translation other than the final send is on
\texttt{src}, \texttt{acc} or \texttt{ans}, which the translation restricts;
these are internal steps. Lookahead is internal by definition. By
Obligation~\ref{obl:fragment} each contribution exists and is unique; by
commutativity and associativity of the aggregator, the value accumulated is
independent of the order of the folds; the count guarantees exactly one
publication. A finite burst of internal steps with a unique end-point followed
by one visible output is weakly bisimilar to that output.
\end{proof}

So \fe{} is synchronous in the strongest available sense: it \emph{is} one step
up to weak bisimulation. Termination comes from \(\mathcal{T}\) (and from
funding, \S\ref{sec:cost}); determinism from \(\mathcal{T}\), from
Obligation~\ref{obl:confluence} for the monad's rewrites, and from the
aggregator's laws. The slogan --- concurrent internally, sequential externally
--- is the content of the proposition.

%=============================================================================
\section{Aggregation}
\label{sec:agg}

\subsection{Parallel composition is the free commutative monoid}
\label{sec:free}

\texttt{PPar} \emph{is} \texttt{HashBag(Proc)}, and a finite bag is the free
commutative monoid on its elements. So for every commutative monoid
\((A, \oplus, \varepsilon)\) and every assignment \(h : \Proc \to |A|\) there is
exactly one monoid homomorphism \(\hat h\) from parallel composition to \(A\),
\[
\hat h(P \mid Q) = \hat h(P) \oplus \hat h(Q),
\qquad
\hat h(\texttt{0}) = \varepsilon .
\]
Programmable aggregation is \emph{defined} by this universal property. The
default rule of \S\ref{sec:rule} is the meaning of \fe{}, and every other
aggregator is post-composition with a homomorphism out of it:
\[
\texttt{foreach(\dots) yield( }r\texttt{ ) by }A\ P
\quad\text{means}\quad
r!\big(\hat h\,(\Pi\,\mathsf{contrib})\big).
\]
The assignment \(h\) needs no syntax. The body is typed in the carrier, so the
body \emph{is} \(h\): it maps each binding to an element of \(|A|\).

Three consequences follow, and they are the reason to choose this path.
\begin{itemize}[nosep]
\item \textbf{No new semantics.} Parallel composition stays the reference
meaning; aggregators add structure on the output side only.
\item \textbf{Fusion is a theorem.} An implementation folds each contribution
into the accumulator as it arrives instead of building the parallel composition
and then folding it. That is correct exactly because \(\hat h\) is a
homomorphism, and the specification never has to mention it.
\item \textbf{Order stays invisible.} The argument that reordering by the
concurrency is undetectable --- made first for parallel composition --- holds for
every commutative aggregator, for the same reason.
\end{itemize}

\subsection{Aggregators are theory instances}
\label{sec:instances}

The \texttt{by} clause names a \emph{theory instance}: an elaborated theory
expression --- an application or a \texttt{let}-bound name, of the grammar
theory bodies already use --- whose presentation descends from
\texttt{CommutativeMonoid}. Integer summation is one replacement away from the
universal-algebra module:
\begin{lstlisting}[language=mettail]
Module Arith {
  import "UnivAlg.module" as u
  Theory IntSum() {
    let s  = u.EmptySet() in (
    let m  = u.Monoid(s) in (
    let cm = u.CommutativeMonoid(m) in (
    cm
      Exports      { Elem => Int; }
      Replacements { Zero => IZero . |- "0" : Int;
                     Plus => IAdd  . x:Int, y:Int |- x "+" y : Int; }
      Rewrites     { IAddEval : (IAdd n k) ~> [n + k];   // host integer primitive
    })))
  }
}
\end{lstlisting}
and in a program,
\begin{lstlisting}
import "Arith.module" as a

foreach( i <- { 1, 3, 5, 7, 9 } where i > 4 ) yield( rslt ) by a.IntSum() i
  ~>  rslt!( 21 )
\end{lstlisting}

\paragraph{The default is an instance too.} Without a \texttt{by} clause, the
aggregator is the \texttt{ParMonoid} instance of the theory that defines
\texttt{PPar} in the program's own language. Nothing about the default is
special-cased; it is the instance that happens to be free.

\paragraph{Commutativity is checked by origin, not by name.} Instances were
chosen over named operations because of what the elaborator can then decide.
Every element of an elaborated presentation carries its origin token, and a
replacement keeps the origin of what it renames. So the elaborator can check
that the instance's operation \emph{is} the \texttt{Plus} introduced by
\texttt{u.CommutativeMonoid}, carrying that theory's associativity, unit and
commutativity equations by origin, and not merely an operation called
\texttt{Plus}. This is the mechanism that already distinguishes a shared
\texttt{Name} from two categories that happen to share a spelling. Two
diagnostics follow.
\begin{itemize}[nosep]
\item The diagnostic \texttt{[aggregator-not-commutative]} is raised when the
instance's operation does not descend from \texttt{CommutativeMonoid}'s \texttt{Plus}. An operation that is
associative but not known to be commutative is an ordered aggregator
(\S\ref{sec:stages}) or nothing.
\item \texttt{[aggregator-ambiguous]}: the instance carries more than one
commutative-monoid structure on its carrier --- a rig carries two, by
construction. The clause must then select one by the parameter that supplied it,
\texttt{by r.add}.
\end{itemize}

\paragraph{Elaboration is compile time.} Instances are elaborated by the client
toolchain, as all theory expressions are. An aggregator is therefore fixed by
the program text; it cannot be chosen by data received at run time. That is the
right default for a construct whose cost must be bounded before admission.

\paragraph{What remains an obligation.} Origins establish that the operation is
\emph{declared} associative, commutative and unital. That the rewrites which
\emph{compute} it respect those equations --- that \texttt{IAddEval} agrees with
commutativity, say --- is a confluence obligation on the instance, of the same
kind as Obligation~\ref{obl:confluence}, and the operation's own evaluation must
lie in \(\mathcal{T}\).

\subsection{Tree-shaped accumulation}
\label{sec:tree}

With a commutative monoid, the accumulator need not be a single serial
channel. Let every contribution be a singleton partial result, paired with a
count, and let partial results combine wherever they meet:
\begin{lstlisting}
new src, acc in {
  [[ m ]]_src
| for( @x <= src ){
    match x {
      ptrn where cond => acc!( 1, [[ P ]] )[*]
      _               => acc!( 1, eps )
    }
  }
| for( @(k1, a1) <= acc & @(k2, a2) <= acc where k1 + k2 <  n ){ acc!(  k1 + k2, a1 (+) a2 )[*] }
| for( @(k1, a1) <- acc & @(k2, a2) <- acc where k1 + k2 == n ){ rslt!( a1 (+) a2 )[*] }
}
\end{lstlisting}
with \texttt{(+)} and \texttt{eps} the aggregator's operation and unit, and
two further rules for \(n \le 1\): an empty source publishes \texttt{eps}, and a
single contribution is published as it stands. Three observations.
\begin{itemize}[nosep]
\item The pair of a count and a value lives in the product
\((\mathbb{N},+,0)\times(A,\oplus,\varepsilon)\), itself a commutative monoid.
The completion test is a second aggregator running beside the first, and the
deficit of \S\ref{sec:refsem} for dependent groups is the same construction over
\(\mathbb{Z}\).
\item Combination is no longer a chain of length \(n\); partial results merge
pairwise wherever they meet, so the depth can be logarithmic. This is the
combiner of MapReduce~\cite{DeanGhemawat04}, and it is the shape of parallelism
the GPU deployment of the node wants.
\item Accumulation is itself a rewrite system on a bag of partial results, in
the style of the communication rule, and it publishes exactly once.
\end{itemize}
The serial translation of \S\ref{sec:refsem} and this one are weakly bisimilar
by commutativity and associativity, and either may serve as the reference.

\subsection{The staged path}
\label{sec:stages}

Each further stage adds one algebraic requirement on the aggregator, checked by
origin in the same way, and costed by the same per-combine rule.

\begin{center}
\small
\begin{tabular}{@{}L{0.12\textwidth}L{0.24\textwidth}L{0.28\textwidth}L{0.28\textwidth}@{}}
\toprule
\textbf{Stage} & \textbf{Aggregator descends from} & \textbf{Examples} &
\textbf{Requirement on the source} \\
\midrule
1 & \texttt{CommutativeMonoid} & parallel composition, sums, counts, products &
\texttt{toBag} \\[3pt]
2 & a \texttt{Semilattice}: commutative, idempotent & union, maximum, boolean
or & \texttt{toBag} \\[3pt]
3 & \texttt{Monoid} only & lists, strings, function composition &
\texttt{toBag} and \texttt{ordered}; \texttt{\&} needs \texttt{commutative} \\[3pt]
4 & a \(\mathcal V\)-\texttt{Semimodule} & weighted bags, expectations,
amplitudes & \texttt{toBag}; a graded \texttt{where} \\
\bottomrule
\end{tabular}
\end{center}

\paragraph{Stage 2: idempotent aggregators.} A semilattice is a commutative
monoid with \(a \oplus a = a\). It needs a theory that the universal-algebra
module does not yet contain: a \texttt{Semilattice}, taking a commutative
monoid as its parameter and adding the idempotence equation. Its first use settles a question the earlier draft left
open: whether a \texttt{Set} source should deduplicate. It need not. The
comprehension visits binding instances (\S\ref{sec:ddl}), and
\begin{lstlisting}
foreach( x <- { 1, 2 } ; y <- { 3 } ) yield( r ) y                 ~>  r!( 3 | 3 )
foreach( x <- { 1, 2 } ; y <- { 3 } ) yield( r ) by u.SetUnion() y  ~>  r!( Set(3) )
\end{lstlisting}
Deduplication is the choice of an idempotent aggregator. Semilattices are also
the coalition-logic container of the generated logics~\cite{F1R3Dev}, so this
stage is where aggregation and observation meet.

\paragraph{Stage 3: ordered aggregators.} A monoid that is not commutative ---
list concatenation, string building, composition of functions --- breaks the
argument that order is invisible. Contributions are then tagged with their
position in the source's \texttt{ordered} enumeration and collected in the
monoid of finite maps from positions under disjoint union, which \emph{is}
commutative; the non-commutative fold is applied once, in index order, at
publication. Tree-shaped accumulation is preserved. The source must declare
\texttt{ordered}: \texttt{Seq} does, and \texttt{Bag} and \texttt{Set} have no
canonical order, so an ordered aggregator over them is refused with
\texttt{[aggregator-needs-order]}. This stage is also where Scala's
\texttt{for}/\texttt{yield} lives: yielding into the source monad itself, with
\(\unit\) and concatenation, preserves its order.

\paragraph{Stage 4: semimodules.} A graded \texttt{where} clause takes values in
a resolution algebra \(\mathcal V\)~\cite{GradedWhere}. Its weight cannot be
applied to a contribution unless the aggregator admits scalars: a
\(\mathcal V\)-semimodule, a commutative monoid with an action of \(\mathcal V\)
distributing over \(\oplus\). The graded contribution is \(w \cdot
\nf(P\sigma)\). The free \(\mathcal V\)-semimodule on processes is the weighted
bag, in exactly the sense that parallel composition is the free commutative
monoid, so Stage~4 is the universal-property construction of
\S\ref{sec:free} one level up. Section~\ref{sec:graded} returns to it.

%=============================================================================
\section{Sequencing}
\label{sec:seq}

Rholang~1.1's synchronous send sequences on an acknowledgement that the
receiver may or may not send. \fe{} has an acknowledgement it always sends: its
result. So it can stand on the left of \texttt{;}, and the continuation waits on
the result channel:
\begin{lstlisting}
foreach( p <- m where c ) yield( r ) P ; K
  ==  new r in { [[ foreach( p <- m where c ) yield( r ) P ]] | for( @_ <- r ){ K } }
\end{lstlisting}
A continuation that needs the value binds it instead,
\begin{lstlisting}
foreach( p <- m where c ) yield( r ) P | for( @ys <- r ){ K }
\end{lstlisting}
so no separate value form of the comprehension is needed: the acknowledgement
and the result are the same datum. In the first form \texttt{r} is bound by the
sugar, and is private to the comprehension and its continuation.

\paragraph{Why \fe{} has the right.} A process may appear on the left of
\texttt{;} only if the continuation is guaranteed to be released. The
synchronous send cannot offer that guarantee: it depends on a convention
honoured by the receiver. \fe{} can. Its guards and bodies terminate because
they are well typed in System F (later, in linear logic) and the translation
reflects divergence (Obligation~\ref{obl:fragment}); its count reaches the
measure; and so its publication is certain whenever it is funded. When it is
not funded, nothing is published, \texttt{K} is not released, and the deployer
is notified. Termination is what earns the place on the left of \texttt{;}, and
it is the type system that grants it; funding decides only whether this run
gets there.

\paragraph{Chains.} Mixed chains need no further machinery. The chain
\begin{lstlisting}
x!?(Q); foreach( y <- m ) yield( r ) P; z!?(R); K
\end{lstlisting}
is nested \texttt{new}/\texttt{for} pairs, and \texttt{;} means ``then''
throughout: between receipt groups, after a synchronous send, and after a
\fe{}. The earlier draft's sequenceable fragment of bodies is no longer needed:
completion is defined by publication, for every comprehension.

%=============================================================================
\section{Cost}
\label{sec:cost}

\subsection{Two gates}

The cost-accounted form of the comprehension is
\begin{lstlisting}
{% foreach( ptrn <- mndInst where cond ) yield( rslt ) by A {% P %}[t] %}[s]
\end{lstlisting}
and its two gates fund different things.

\paragraph{The outer gate is per binding.} In the core, a gate funds one event,
the communication. In \fe{}, the event the comprehension witnesses is each
\emph{binding}, so \texttt{[s]} funds one token per candidate --- every element
of \(\mathsf{toBag}(m)\), whether or not its pattern and guard succeed --- plus
the combinations of the accumulator and the one publication. The principle is
the one the developer document states, the witness is where cost attaches, with
a different multiplicity.

\paragraph{The inner gate funds lookahead.} \verb|{% P %}[t]| funds the
reduction of a surviving body to normal form. Because lookahead runs in
isolation, the inner gate is spent during the comprehension, not by a later
consumer, and its budget depth bounds the work: it is the terminating fragment
of \S\ref{sec:lookahead} made concrete. Lookahead of the guard and of the
aggregator's operation is charged in the same way, on the lane the program
assigns.

\paragraph{Per-group gates.} The per-clause signed binds of the core carry over:
\begin{lstlisting}
{% foreach( {% x <- m %}[s1] ; {% y <- n(x) %}[s2] ) yield( r ) {% P %}[t] %}[s0]
\end{lstlisting}
funds the bindings of each group from its own lane. This matters for dependent
groups, whose second-group candidates may be supplied by a different party.

\paragraph{The syntax is forced.} In the native four-sort grammar, continuation
and payload slots have the signed-term sort \(\T\). Giving \fe{} the same
treatment --- the comprehension in \(\T\), the body in \(\T\) --- produces
exactly the surface above. An unwrapped \fe{} cannot be written in the
cost-accounted surface, and wrapping by construction continues to hold. The
continuation after \texttt{;} is signed too, exactly as the continuation inside
a desugared \verb|!?| is:
\begin{lstlisting}
{% foreach( ptrn <- mndInst ) yield( r ) {% P %}[t] %}[s] ; {% K %}[u]
\end{lstlisting}

\begin{proposition}[The outer cost does not depend on the filter]
\label{prop:filterindep}
The number of tokens a \fe{} draws on its outer gate for bindings equals
\(\meas(m)\), whether or not patterns and guards succeed. The number of body
lookaheads it funds on its inner gate is the number of survivors, and is bounded
by the same count.
\end{proposition}

\begin{proof}
Every element of \(\mathsf{toBag}(m)\) is matched and guarded, and contributes
--- its body's normal form or the unit --- so every element passes the outer
gate. Only survivors are looked ahead, and survivors are a subset of elements.
\end{proof}

The proposition closes a hole that would otherwise exist. If only the inner
gate charged, a \fe{} over a very large collection with a guard that is always
false would cost almost nothing, and the traversal would be free. Per-binding
charging prices every candidate, and a single estimate bounds both lanes.

\subsection{The atomic outer gate}
\label{sec:atomicgate}

The number of candidates \(n = \meas(m)\) is computable when the gate opens, so
for a single group the outer gate can be an \(n\)-ary join on the purse: one
communication that takes everything the binding lane needs or takes nothing.
\begin{itemize}[nosep]
\item If the purse holds enough, the join fires, the comprehension runs as a
burst of internal steps, and it ends in its one publication.
\item If it does not, nothing is consumed, and the comprehension is a non-event
--- exactly as an unfunded communication is in the core.
\end{itemize}
Because bodies run in isolation, the comprehension has no effects other than
its publication, so on the outer lane it is all-or-nothing in the strongest
sense.

The inner lane is different, and deliberately so. Typing guarantees that a
body's lookahead terminates, but not within any particular budget. A body whose
lookahead exhausts its \texttt{t} budget never contributes, the count never
reaches \(n\), and nothing is published. The outer tokens are spent, nothing is
reserved or refunded, and the deployer is notified: the node already records
allocation denials in the deploy's budget trace, and an exhausted lookahead is
one more such denial. To consumers the comprehension is a non-event. That is
the policy, and it is sufficient: a comprehension that is not sufficiently
funded does not complete. Dependent groups reopen the outer question, since
their total is known only as the first group expands
(Question~\ref{q:dependent}).

\subsection{Precharge estimates}
\label{sec:precharge}

Per-binding charging moves the variability of \fe{}'s cost out of the semantics
and into estimation. No static guarantee of completion is required, since an
underfunded comprehension simply does not complete. What an estimate buys is
the deployer's ability to fund correctly the first time, and on the outer lane
it can be exact once \(\meas(m)\) is bounded.

\begin{remark}[Two uses of ``precharge'']
In the node, \emph{precharge} also names the escrow mechanism --- precharge,
run, refund --- that the cost-accounting work removed in favour of
gate-funded execution. That mechanism is not being reintroduced. Here
\emph{precharge} means the estimate a deployer funds against, and nothing is
escrowed: the \(n\)-ary gate takes exactly what it needs when it
fires.
\end{remark}

\paragraph{What is estimated.} For a single group with \(n = \meas(m)\), an
aggregator whose operation costs at most \(c_\oplus\) to look ahead, and a body
budget of depth \(d_t\),
\[
\text{outer} \;\le\; n \;+\; (n-1)\,(1 + c_\oplus) \;+\; 1,
\qquad
\text{inner} \;\le\; n \cdot d_t ,
\]
counting one token per binding, one per combination plus its operation, and
one for the publication. For dependent groups, \(n\) is
\(\meas(m) + \sum_x \meas(n(x))\). The outer figure is exact given \(n\) and
\(c_\oplus\). The inner figure is not derived from anything: \(d_t\) is the
budget the deployer chooses to give each body, since System F typing supplies
no bound to derive it from. A \fe{} nested inside a body is funded from the
inner lane and draws on that same \(d_t\).

\paragraph{Where bounds come from.} In decreasing order of precision:
\begin{enumerate}[nosep]
\item \emph{A closed instance at preparation time.} The measure is computed and
the bound is exact.
\item \emph{An instance bound by an earlier receipt or match.} Theory types say
nothing about size, and a pattern such as \verb|[a, b, ...rest]| supplies only
a lower bound. The guard on that binding supplies the upper bound:
\begin{lstlisting}
for( @m <- ch where m.size() <= 64 ) {
  {% foreach( x <- m ) yield( r ) {% P %}[t] %}[s]
}
\end{lstlisting}
Here the \texttt{where} clause does double duty: it is the developer's filter,
and it is the premise the estimator reads to produce an exact outer figure. A
deployer who wants that figure before running writes such a guard.
\item \emph{At the gate.} The exact measure is available when the outer gate
opens, and the \(n\)-ary gate uses it. That makes the gate exact, but tells the
deployer nothing in advance.
\end{enumerate}

\paragraph{Two admissibility classes.} A monad whose \(\mathsf{toBag}\) is not
finite on its instances --- lazy streams, say --- cannot be measured, and so can
never be precharged, and its comprehension can never complete. Since the
completion condition of \S\ref{sec:refsem} is the measure, such monads are not
merely hard to estimate but unusable with \fe{}; the elaborator reports
\texttt{[monad-unmeasured]}. Among finite monads, an instance with a known bound
gives the deployer an exact outer estimate, and any other instance is decided at
the gate.

\paragraph{Lanes.} The \texttt{s} and \texttt{t} gates draw on different
per-signature purses, so the proof splits by signature, and the node's disjoint
per-signature lane pools already enforce the split; so does the
notification, which names the lane that ran dry. If \texttt{t} may mention
variables bound by the pattern --- \verb|{% P %}[#x]| --- then each element pays
for its own processing, which is useful when the elements of a collection come
from different signers; the inner lanes then become data-dependent, and an
estimate must range over the elements' signatures rather than over a fixed lane
(Question~\ref{q:hash}).

%=============================================================================
\section{Implementation}
\label{sec:impl}

\subsection{The translation is the specification, not the strategy}

\begin{proposition}[Conservativity]
\label{prop:conservative}
A program that does not use \fe{} has the same reductions, the same costs and
the same observable behaviour with the extension as without it.
\end{proposition}

\begin{proof}[Sketch]
Every new rule has a \fe{} constructor on its left-hand side, and no existing
rule produces one. The \texttt{Monads} builder adds rules only to theories that
declare it. The cost rule attaches only to the new gate.
\end{proof}

Nothing therefore obliges an implementation to compile \fe{} into
\texttt{for}. The translations of \S\ref{sec:refsem} and \S\ref{sec:tree} are
retained as the reference semantics, and an implementation may evaluate \fe{}
natively, as the single rule of \S\ref{sec:rule}.

\begin{obligation}[Graded refinement]
\label{obl:graded}
The native evaluation of every \fe{} is weakly bisimilar to its reference
translation under the graded, signature-labelled transition relation.
\end{obligation}

The native charge is designed to equal the reference translation's metered
communications: one per element on \texttt{src}, one per combination on
\texttt{acc}, and one for the publication. Graded bisimilarity then states both
requirements at once, the same behaviour and the same bill. The relation
already exists in the node's Rocq development, as the graded transition system
and its Hennessy--Milner adequacy results~\cite{f1r3nodeOTE}. Conservativity
itself still needs checking against the reducer, an ordinary compatibility
review.

\subsection{What native evaluation reuses}

Over a built-in collection, the native rule is an iterated match: the spatial
matcher, the binders, and the \texttt{where} guards on match cases with
fall-through, which are already implemented, supply the per-element
pattern-and-guard step. What is new is lookahead --- evaluating a body to normal
form in a private sub-configuration that shares nothing with the tuple space ---
and the fold into the aggregator, which is the instance's own rewrite engine
applied to its operation.

\subsection{What native evaluation must not inherit}

\paragraph{The free-reduction default.} The node's consensus ledger counts
only communication events; non-communication structural reductions ---
\texttt{new}, \texttt{match}, \texttt{if} --- emit diagnostic events that are
not charged. A natively evaluated \fe{}, being match-class work plus lookahead,
would by default cost nothing however large the instance, and with host work
unmetered by design, that would be free unbounded work. The per-binding charge
and the lookahead budget of \S\ref{sec:cost} are therefore mandatory. Both are
charges on language-level events, not on host work, so they remain consistent
with the decision to cap rather than meter host work. A hard protocol cap on
measure sits beside them.

\paragraph{The static demand analysis.} The funding analysis counts the
communication nodes reachable in a normalised term, relies on receiving a term
already desugared, and keeps an explicit seam for a front-end that delivers
something else. A native \fe{} is that something else: it is not a
communication node, and its contribution is data-dependent. The analysis gains
a \fe{} rule, the bounds of \S\ref{sec:precharge}.

\subsection{What changes in the representation}

The built-in collection methods (\texttt{nth}, \texttt{union},
\texttt{toList}, and so on) are first-order, and the method-call node cannot
carry a binder. A native \fe{} needs a new process node carrying the groups'
binders --- most naturally reusing the receipt's bind-and-pattern encoding ---
the result channel, the body, and a reference to the aggregator instance by the
digest of its elaborated presentation. That is a change to the protobuf model,
and so it is consensus-visible. It should be gated on the protocol version, as
the new cost-record fields are, and it must pass the compatibility review the
frontend admission contract requires, since that contract freezes the model and
protobuf files~\cite{f1r3nodeOTE}.

\subsection{What gets simpler}

\begin{itemize}[nosep]
\item \textbf{Replay.} A native \fe{} touches the tuple space once, to publish.
Lookahead and accumulation leave nothing in the event log, so replay
recomputes them, and recomputation is deterministic by
Obligation~\ref{obl:fragment}, Obligation~\ref{obl:confluence} and the
aggregator's laws.
\item \textbf{Sequencing.} \texttt{K} after \texttt{;} is released by the
publication. Natively, the sugar's private receipt can be fused away
altogether.
\item \textbf{Order-then-execute.} The internal steps race with nothing, across
deploys or within one, so a \fe{} introduces no conflicts to mark by signature
other than its publication, which is an ordinary send. It is order-independent
and replay-deterministic by construction.
\end{itemize}

\subsection{Built-in and declared monads}

Two evaluation paths coexist.
\begin{itemize}[nosep]
\item \emph{Built-in collections} get a fast path in the reducer.
\item \emph{Monads declared in a theory} run their \(\mathsf{toBag}\) and
\(\bind\) rules through the installed language's rewrite engine --- the
semantic-kernel service through which installed languages reduce terms. That
work sits today under host-side semantic-work limits rather than the token
ledger. The per-binding charge and the lookahead budget must cover this path
too, or user-declared monads become the free route the built-ins were closed
against. The same holds for aggregator operations defined by a theory's
rewrites.
\end{itemize}

\subsection{Where the frontend stands}

On the integration branch, \texttt{Module} and \texttt{Theory} declarations are
parsed structurally, the existing lowerer emits the node's process
representation directly, and public activation is off~\cite{f1r3nodeOTE}. The
work this note adds lands in four places: the elaborator (the \texttt{Monads}
builder, \texttt{toBag}, its obligations, and the origin check on aggregator
instances with its diagnostics); the type checker (source sorts, the carrier of
the aggregator, System F typing of guards and bodies) and the translation
\(\sem{-}\) that compiles them; the reducer (the native node,
lookahead in isolation, and the charge); and the funding analysis (the demand
rule).

%=============================================================================
\section{Interactions with the rest of \flang{}}
\label{sec:interactions}

\subsection{Channel types}

A channel's type is an elaborated theory, so a \texttt{Monads} declaration
travels with it: a receiver may \fe{} over a received term only when the
channel's theory carries a monad on the term's sort. Symmetrically, the result
channel's type must export the aggregator's carrier, so a consumer of
\texttt{rslt} knows what it receives --- parallel compositions under the default,
integers under \texttt{IntSum}. Two agents sharing a language definition share
its iteration and aggregation structure as well as its grammar.

\subsection{Generated logics}
\label{sec:logic}

The new constructors and rules are visible to the logic generated from the
theory: \(\unit\), \(\bind\) and \(\mathsf{toBag}\) to the spatial connectives,
their rewrites to the behavioural modality. The aggregator's carrier is a sort
like any other, so a consumer's guard on the result is an ordinary formula over
it.

There is a richer option. The logical counterpart of a monad is a \emph{lax
modality}~\cite{FairtloughMendler97, BentonBiermanDePaiva98}, written
\(\lax\varphi\): ``after binding in \(\M\), \(\varphi\).'' A \texttt{Monads}
declaration could generate \(\lax_\M\) in the generated logic, giving
\texttt{where} clauses a way to speak about what a monadic value will deliver
without enumerating it. The lax modality is a type-system-flavoured connective
and raises no new adequacy question: it is consistent with the established
position that Hypercube-generated logics are type systems and are not expected
to be adequate.

\subsection{Graded \texttt{where} clauses}
\label{sec:graded}

Here the choose-versus-sum distinction of \S\ref{sec:intro} does its work.
\begin{itemize}[nosep]
\item In \texttt{for}, a graded guard \emph{weights candidates in a race}: the
weight changes which candidate is chosen.
\item In \fe{}, there is no race. A graded guard must \emph{scale
contributions}: the weight multiplies what a candidate adds to the aggregate.
\end{itemize}
Scaling contributions requires the aggregator to be a semimodule over the value
algebra --- Stage~4 of \S\ref{sec:stages}. Parallel composition has no scalars,
so a graded \texttt{where} under the default aggregator is refused, and the
requirement falls on the aggregator, where it belongs, rather than on the
monad.

\subsection{Simulation}

A \fe{} whose aggregator is the weighted bag, and whose guard is graded, is a
stochastic simulation step: programming in the large, written with the same
surface as programming in the small~\cite{WeightedGSLT}. The ``sum of all
possible futures'' of the equational route to a quantum reading of \rhoc{} is
what such a comprehension publishes.

\subsection{The quantum fence}

With amplitudes as the value algebra, aggregation is a sum of weighted vectors,
and Kleisli composition in the vector-space monad is matrix multiplication ---
which is why the rho--linear-algebra correspondence and the comprehension meet.
But a body that uses its bound variable twice is not linear, and the quantum
reading of graded \texttt{where} clauses is restricted to the linear fragment of
\rhoc{}. The linearity check of that fragment applies to \fe{} bodies unchanged.
The second terminating fragment, Abramsky's proof expressions, is linear by
construction, so it is the natural home for bodies under this reading.

\subsection{Summary}

\begin{center}
\small
\begin{tabular}{@{}L{0.22\textwidth}L{0.70\textwidth}@{}}
\toprule
\textbf{Feature} & \textbf{Interaction with \fe{}} \\
\midrule
Theory-level monads & laws hold modulo erasure in \(\cost(T)\), \(\hist(T)\)
(Prop.~\ref{prop:modulo-erasure}) \\[2pt]
Equations & laws are obligations, never equations; no higher-order matching \\[2pt]
Rewrites & interpretation and \texttt{toBag} are rewrites; confluence
obligation \\[2pt]
Algebra of theories & origin tokens; \texttt{[join-collision]},
\texttt{[dangling-monad]}; arity drift shared with \texttt{PPar} \\[2pt]
Universal-algebra tower & aggregators are \texttt{CommutativeMonoid}
instances, checked by origin; default \texttt{ParMonoid} \\[2pt]
Arrows & only \texttt{<-} on sources; consumers use all three on the
result \\[2pt]
\texttt{;} and \texttt{\&} & dependent bind via \texttt{toBag}; \texttt{\&}
free under commutative aggregators \\[2pt]
\texttt{where} & rejected candidates contribute the aggregator's unit \\[2pt]
Synchronous send & \fe{} sequences unconditionally; \verb|!?| by
convention \\[2pt]
Cost & two gates; per-binding outer charge; lookahead budget; precharge by
measure \\[2pt]
Channel types & monad travels with the source's theory; result typed by the
carrier \\[2pt]
Generated logics & new constructors visible; optional lax modality \\[2pt]
Graded \texttt{where} & needs a semimodule aggregator (Stage~4) \\[2pt]
Quantum reading & linear bodies only \\[2pt]
Order-then-execute & one publication; replay-deterministic \\
\bottomrule
\end{tabular}
\end{center}

%=============================================================================
\section{A worked example: the parametric \texorpdfstring{\rhoc{}}{rho} calculus}
\label{sec:rho}

The prelude monads of \S\ref{sec:ddl} are collections. The example that
motivates Step~2 is of a different kind, and it is the calculus itself. Two
working notes treat the syntax of \rhoc{} as a type constructor: \emph{Where
Substitution Happens} parameterises it by a set \(D\) of data carried on names
and messages and a set \(V\) of variables~\cite{WSH}, and \emph{Oriented
Processes} specialises \(D\) to a group \(G\) and lets \(G\) act on
processes~\cite{Oriented}. This section expresses that construction in the
extensions proposed here. It is a monad, and the group action is a
distributive law.

\subsection{The construction}

In the paired-message calculus of~\cite{WSH}, messages are pairs
\(\langle P, d\rangle\) of a code and a datum \(d\in D\), names are quoted
messages, and variables form their own syntactic category, occurring directly
only under the unary drop \(\ast v\). Communication substitutes a received name
for a bound variable, and substitution is \emph{semantic}: it commutes with
every constructor, quotation included, and its one non-trivial case is
\[
(\ast v)\{@\langle p, d\rangle / v\} \;=\; p .
\]
Separating variables from names is what makes \(V\) a parameter: the binding
discipline lives in one place and can be chosen independently of the name
theory~\cite{WSH}.

\emph{Oriented Processes} works with the pattern-matching variant, in which a
receipt binds both the code and the datum of what it receives, and fixes
\(D=G\). Variables occur \emph{decorated}, \([g]X\) and \([g]\delta\),
carrying the action pending on them until they are instantiated. The action
\(g\cdot-\) multiplies every literal datum on the left at every depth of
quotation and is recorded, not applied, at variables:
\(g\cdot[k]X=[gk]X\). Substitution discharges the record,
\(([g]X)\sigma = g\cdot\sigma(X)\). On closed terms the action is the image of
left multiplication \(L_g\) under the functor that the parametric syntax
defines~\cite{Oriented}.

\subsection{The syntax monad}

Write \(\mathsf{Rho}_D(V)\) for the processes over data \(D\) whose free
variables lie in \(V\). This is a monad in \(V\) for each fixed \(D\), and it is
the standard monad of syntax with variables and binding: the term is a tree
whose leaves are variables, and the monad's bind is substitution at the
leaves~\cite{BirdPaterson99, AltenkirchReus99, FiorePlotkinTuri99}.
\begin{itemize}[nosep]
\item \(\unit(v) = \ast v\): a variable, as the process that runs it.
\item \(\bind(P, \sigma)\), for \(\sigma : V \to \mathsf{Rho}_D(W)\), replaces
every \(\ast v\) in \(P\) by \(\sigma(v)\), under quotation and avoiding
capture. A receipt's body lives in \(\mathsf{Rho}_D(V+1)\), and bind extends
\(\sigma\) by the identity on the bound variable as it passes the binder.
\end{itemize}
The left unit law is the base case of the definition, the right unit law says
that substituting each variable for itself is the identity, and associativity is
the substitution lemma. All three are properties of capture-avoiding
substitution, not of any particular theory.

That last observation shapes the surface. For a syntax monad, the
\texttt{Monads} builder does not need the interpretation written out by cases.
The toolchain already generates capture-avoiding substitution from the binders
declared in \texttt{Terms}~\cite{F1R3Dev}, so it suffices to say which
constructor is the variable former:
\begin{lstlisting}[language=mettail]
Theory ParamRho() {
  Types {
    Proc(D, V);   Name(D, V);   Msg(D, V);
  }
  Terms {
    PZero  . |- "0" : Proc(D,V);
    PDrop  . v:V |- "*" v : Proc(D,V);
    POut   . x:Name(D,V), m:Msg(D,V) |- x "!" "(" m ")" : Proc(D,V);
    PIn    . x:Name(D,V), ^v.p:[V -> Proc(D,V)]
           |- "for" "(" v "<-" x ")" p : Proc(D,V);
    PPar   . ps:HashBag(Proc(D,V)) |- "{" ps.*sep("|") "}" : Proc(D,V);
    PBDrop . x:Name(D,V), u:Name(D,V) |- "*" "(" x "," u ")" : Proc(D,V);
    MPair  . p:Proc(D,V), d:D |- "<" p "," d ">" : Msg(D,V);
    NQuote . m:Msg(D,V) |- "@" m : Name(D,V);
  }
  Equations {
    (PBDrop (NQuote m) u) == (POut u m);        // E2: the binary drop
  }
  Rewrites {
    RComm : (PPar {(PIn x ^v.p), (POut x m), ...rest})
              ~> (PPar {(subst ^v.p (MCode m)), ...rest});
    RPar  : if S ~> T then (PPar {S, ...rest}) ~> (PPar {T, ...rest});
  }
  Monads {
    Rho over Proc(D, _) by substitution at PDrop;
    toBag : Occurrences;   ordered : Positions;
    laws : generated;
  }
}
\end{lstlisting}
Three features of the declaration are new, and each is a small, definite
addition to the proposal.

\paragraph{\texttt{over Proc(D, \_)}.} The monad is in one parameter; the other
is fixed. This is the general shape of a parametric type constructor in the
DDL, and the hole marks the monadic position.

\paragraph{\texttt{by substitution at PDrop}.} A syntax monad is declared by
naming its variable former. \(\bind\) is the generated substitution, specialised
so that the variable former is the leaf: a substitution
\(\sigma : V\to\mathsf{Rho}_D(W)\) replaces each \texttt{(PDrop v)} by
\(\sigma(v)\). This is exactly semantic substitution, because variables occur
directly only under the drop.

\paragraph{\texttt{laws : generated}.} The laws of a syntax monad follow from
the generated substitution's own correctness, so they are discharged once, for
the toolchain, and not once per theory. This is a better position than the
prelude collections are in, whose laws are discharged monad by monad.

\begin{remark}[The communication rule is a bind]
\label{rem:commbind}
In \texttt{RComm} the reduct \texttt{(subst \^{}v.p (MCode m))} is
\(\bind(p, \sigma)\) for the Kleisli arrow that sends the bound variable to the
code of the received message --- \texttt{MCode} is the projection taking
\texttt{(MPair p d)} to \texttt{p} --- and every other variable to itself. So every
communication performs one bind in the syntax monad: the space monad of
\S\ref{sec:intro} delivers its datum by binding the continuation in the data
monad of the calculus's own syntax. The two comprehensions meet exactly at the
witness. The envelope erasure of~\cite{WSH} --- a receiver can never record the
datum of the message it receives --- has a one-line reading: the Kleisli arrow
of \texttt{RComm} factors through the projection
\(\mathsf{Msg}_D(V)\to\mathsf{Rho}_D(V)\) that forgets \(D\). The
pattern-matching binder of~\cite{Oriented}, which binds code and datum
together, replaces that arrow with one into a two-sorted substitution that
forgets nothing.
\end{remark}

\begin{remark}[Cost]
The per-binding charge of \S\ref{sec:cost} applies to a user-level \fe{} over
the syntax monad. It is not proposed for the bind inside \texttt{RComm}, which
remains part of a single witnessed communication and is charged as one, as it
is today. Charging that substitution per occurrence would change the cost model
of the core, which this note does not do.
\end{remark}

\subsection{The functor in \texorpdfstring{\(D\)}{D}}

\(D\) occurs only covariantly, in the datum of a message, so
\(\mathsf{Rho}_{(-)}(V)\) is a functor in \(D\) for each \(V\): a map
\(f : D\to D'\) acts on every datum, at every depth of quotation. Step~2 should
generate \(\fmap\) for every covariant sort parameter, in the manner of a
derived functor instance; no declaration is needed. With \(D=G\), the action
of~\cite{Oriented} on closed terms is this generated map at left
multiplication:
\[
g\cdot P \;=\; \fmap_D(L_g)(P) \qquad (P \text{ closed}).
\]
On open terms the generated map is not enough. It moves the literal data but
leaves variables alone, and a variable that is later instantiated would then
escape the action. That is the problem deferral solves, and its solution is a
second monad.

\subsection{Orientation is a distributive law}

Let \(\mathsf{Frame}_G(X) = G\times X\), the writer monad of the group \(G\)
viewed as a monoid. For a group it is the free \(G\)-set monad: \(\unit(x) =
(e, x)\), and binding multiplies frames. Its elements are exactly the decorated
variables of~\cite{Oriented}: \([g]X\) is the pair \((g, X)\), and \(X\) is
\([e]X\). Data are decorated the same way, with a literal datum \(g\) read as
\(g\) applied to the origin, so the data of the oriented calculus are
\(\mathsf{Frame}_G(\Delta + 1)\) for data variables \(\Delta\).

The oriented calculus is the composite
\[
\mathsf{Orient}_G \;=\; \mathsf{Rho}_{\mathsf{Frame}_G(\Delta+1)} \circ \mathsf{Frame}_G ,
\]
processes whose variables carry frames, and the action is the distributive
law that makes the composite a monad~\cite{Beck69}:
\[
\lambda : \mathsf{Frame}_G\circ\mathsf{Rho} \Rightarrow
\mathsf{Rho}\circ\mathsf{Frame}_G,
\qquad
\lambda(g, P) = g\cdot P .
\]
It takes a frame sitting outside a process and pushes it to the leaves:
multiplying data by the generated map, and recording at each variable. The
composite's bind is substitution that discharges the record,
\(([g]X)\sigma = g\cdot\sigma(X)\), which is the substitution
of~\cite{Oriented} verbatim.

\begin{proposition}[The action is a distributive law]
\label{prop:beck}
\(\lambda\) satisfies the four laws of a distributive law of
\(\mathsf{Frame}_G\) over \(\mathsf{Rho}\), and each is a result already proved
in~\cite{Oriented}.
\end{proposition}

\begin{center}
\small
\begin{tabular}{@{}L{0.22\textwidth}L{0.30\textwidth}L{0.38\textwidth}@{}}
\toprule
\textbf{Law} & \textbf{In the oriented calculus} & \textbf{Where it is proved} \\
\midrule
unit of \(\mathsf{Rho}\) & \(g\cdot[e]X = [g]X\) &
the definition of the action, at a variable \\[3pt]
unit of \(\mathsf{Frame}_G\) & \(e\cdot P = P\) &
the automorphism lemma, part (i) \\[3pt]
multiplication of \(\mathsf{Frame}_G\) & \(g\cdot(h\cdot P) = (gh)\cdot P\) &
the automorphism lemma, part (i) \\[3pt]
multiplication of \(\mathsf{Rho}\) & \(g\cdot(P\sigma) = (g\cdot P)\sigma\) &
the lemma on substitution and matching \\
\bottomrule
\end{tabular}
\end{center}

\begin{proof}[Sketch]
The first two laws are the clauses of the definition and the identity case of
the automorphism lemma. The third is its composition case. For the fourth,
both sides agree at a decorated variable, where
\((g\cdot[k]X)\sigma = (gk)\cdot\sigma(X) = g\cdot(([k]X)\sigma)\), and at every
other constructor by induction; this is the calculation of the substitution
lemma, and deferral is what makes \(\sigma\) the same on both sides.
\end{proof}

The proposition also says why deferral was the right design in~\cite{Oriented}.
Without it, the action would be the generated map in \(D\) alone, which commutes
with nothing at the variables, and the composite would not be a monad. With it,
\(\mathsf{Orient}_G\) is a monad, and everything the extensions of this note
offer a monad becomes available to oriented processes.

In the proposed surface the construction is one theory, parameterised by a
group:
\begin{lstlisting}[language=mettail]
Theory OrientedRho(grp: u.Group) {
  ParamRho() \/ grp
    Monads {
      Frame  over Writer(Elem);                  // prelude: X |-> Elem x X
      Orient =  Rho after Frame by generated action;
      toBag : Occurrences;   ordered : Positions;
      laws : generated;
    }
}
\end{lstlisting}
\texttt{Writer} is a prelude monad parameterised by a monoid, here the group's
carrier; it is the first instance of a monad transformer as a theory parameter
(\S\ref{sec:modules}). \texttt{after \dots{} by} composes two declared monads
through a distributive law, and \texttt{generated action} names the law the
toolchain can build: the generated map in \(D\) for data, and the record at the
variable former for variables. Its four obligations are the rows of the table
above, and for \texttt{OrientedRho} they are already discharged. The chaos
constant, idempotent summation and the oriented communication rule
of~\cite{Oriented} are added as further \texttt{Terms}, \texttt{Equations} and
\texttt{Rewrites} in the ordinary way; none of them touches the monad, because
none of them is a variable former.

\subsection{The capabilities}

\begin{center}
\small
\begin{tabular}{@{}L{0.16\textwidth}L{0.36\textwidth}L{0.38\textwidth}@{}}
\toprule
\textbf{Entry} & \(\mathsf{Rho}\) and \(\mathsf{Orient}_G\) &
\textbf{Consequence for \fe{}} \\
\midrule
\texttt{toBag} & the multiset of variable occurrences; a monad morphism, since
the occurrences of \(\bind(P,\sigma)\) are the union, over the occurrences
\(v\) of \(P\), of the occurrences of \(\sigma(v)\) &
\fe{} visits occurrences; the measure counts them, on canonical forms (below) \\[3pt]
\texttt{ordered} & yes: positions in the canonical form enumerate the
occurrences &
ordered aggregators are allowed (Stage~3) \\[3pt]
\texttt{commutative} & no; syntax monads are not commutative, and the
commutativity of \texttt{|} does not make them so &
\texttt{\&} is free under commutative aggregators and refused under ordered
ones \\
\bottomrule
\end{tabular}
\end{center}

\paragraph{The measure needs canonical forms.} Obligation~\ref{obl:tobag}
requires \(\mathsf{toBag}\) to be well defined on the theory's quotient:
congruent terms must have the same bag. For \(\mathsf{Rho}\) the count of
variable occurrences is invariant under structural congruence: the monoid laws
and E2 neither duplicate nor erase variables. For \(\mathsf{Orient}_G\) it is
not. Distributivity, \(P\mid(Q+R)\equiv P\mid Q + P\mid R\), duplicates the
occurrences of \(P\), and chaos, \(\chaos\mid P\equiv\chaos\), erases them. Two
congruent terms then have different bags, and since matching is modulo the
equations, a receiver cannot choose which representative it is sent.
\(\mathsf{toBag}\) must be taken on a canonical form, which the node already
computes as the section of the structural quotient used for
hashing~\cite{f1r3nodeOTE}, and for \(\mathsf{Orient}_G\) that form is the
distributed one: the measure, and so the precharge, is the size of the expanded
term, which can be exponentially larger than the source.

\subsection{\texttt{foreach} over a process}

At run time, \rhoc{} reduces closed processes, and a quoted process received on
a channel has no free variables. A \fe{} over its object variables would have
measure zero. The useful instance is the one in which \(V\) is a data sort of
\emph{holes}: \(\mathsf{Rho}_D(\mathsf{HoleId})\) is the sort of process
templates, closed as far as the calculus is concerned, whose leaves are holes
to be filled, in the manner of the typed holes of the plausible-fiction
marketplace~\cite{PF}.

Filling a template is the monad's bind, and the bind is not a commutative
monoid on contributions, so it is not an aggregator. It is reached in two
steps that keep the comprehension pure. The aggregator collects a
\emph{substitution} --- finite maps from hole occurrences to processes, under
union, a commutative monoid because each occurrence is visited once and every
union is therefore disjoint --- and the consumer applies it:
\begin{lstlisting}
for( @tmpl <- templates where tmpl.holes() <= 32 ) {
  {% foreach( h <- tmpl where Fillable(h) ) yield( fills ) by Subst()
       {% { h : Fill(h) } %}[t] %}[s]
| for( @sub <- fills ){ filled!( tmpl.subst( sub ) ) }
}
\end{lstlisting}
Each surviving occurrence contributes a singleton map whose value is the
lookahead of \texttt{Fill(h)}; a rejected occurrence contributes the empty map,
and its hole stays a hole. The outer gate draws one token per hole occurrence,
the inner one per filled hole, and the guard on the receipt is the bound the
precharge proof reads (\S\ref{sec:precharge}). Under the default aggregator the
same comprehension would publish \texttt{Fill(h1) | Fill(h2) | \dots}: the fills
side by side and detached from the template, which is meaningful, and
different. Over \(\mathsf{Orient}_G\) each hole carries a frame, and
\texttt{Fill(h)} is written in the hole's own coordinates and transported when
the consumer applies the bind: the local frames of~\cite{Oriented}, delivered
through the comprehension.

\subsection{Symmetry reduces the precharge}

The orbit-counting theorem of~\cite{Oriented} has a direct consequence for cost.

\begin{proposition}[Symmetry-reduced precharge]
\label{prop:symprecharge}
Let \(P\) be an oriented template with finitely many hole occurrences, let
\(H=\Stab(P)\), and suppose the body and guard of a \fe{} over \(P\) are
\(H\)-equivariant: \(k([h]x) = h\cdot k(x)\) and the guard holds at \([h]x\)
exactly when it holds at \(x\). Then the comprehension needs only one body
lookahead per \(H\)-orbit of hole occurrences, the number of orbits is
\(\meas(P)/|H|\), and that number is an integer.
\end{proposition}

\begin{proof}[Sketch]
Hole occurrences are atoms in the sense of~\cite{Oriented}, and the action maps
them to hole occurrences, so they form an \(H\)-invariant subset of the free
atoms. The stabiliser acts freely on the free atoms by the orbit-counting
theorem, hence freely on this subset, so every orbit has exactly \(|H|\)
members and \(|H|\) divides their number. By equivariance, the normal form of
the body at \([h]x\) is the action of \(h\) on the normal form at \(x\), which is
applied as a recorded frame rather than computed by a fresh lookahead.
\end{proof}

When the hypotheses hold, the inner lane's precharge is divided by \(|H|\),
and the divisibility theorem is what makes the reduced figure exact rather than
a rounded estimate. The outer lane is unchanged: every occurrence is still a
binding and still contributes, so the reduction applies to the cost of looking
ahead at bodies and guards, not to the traversal or the accumulation. This
is the symmetry reduction of~\cite{Oriented} applied to a comprehension instead
of to a state space: the symmetry is read off the template, not declared.

\subsection{Dynamic frames, and where the word problem lands}

The final extension of~\cite{Oriented} lets decorations be words in received
data, \([w]X\) with \(w\) in the free product of \(G\) with the free group on
the data variables. In the terms of this note, the frame monad changes from
the writer \(G\times(-)\) to the free-group-over-\(G\) monad
\(\Delta\mapsto G * F(\Delta)\), whose bind evaluates words by substituting
words for variables. The composite is formed the same way.

What changes is equality. Matching in \flang{} is modulo the equations of the
theory, and the frame sort is now the carrier of \texttt{u.Group} with variables.
Deciding that two decorations are equal is matching modulo the group axioms in
the \texttt{Equations} block of \texttt{Group}, which is to say the word problem
for \(G * F(\Delta)\). The cost that~\cite{Oriented} identifies is therefore
not a new mechanism. It is the second rung of the DDL --- matching up to
equations --- meeting a theory whose equational problem is as hard as the group's
word problem, and it is paid wherever such terms are matched, including in the
guards of a \fe{}.

%=============================================================================
\section{A first tranche}
\label{sec:tranche}

\begin{enumerate}[nosep]
\item The prelude declares \texttt{Seq}, \texttt{Bag} and the finite powerset on
the three built-in collections, each with \texttt{toBag}, with
\texttt{ordered} on \texttt{Seq} and \texttt{commutative} on the other two.
\item \fe{} with \texttt{<-}, \texttt{;} between groups, \texttt{\&} within
them, crisp \texttt{where}, and \texttt{yield( rslt )}.
\item Lookahead for guards and bodies, with the terminating fragment given by
System F: a System F theory declared in the DDL, guards and bodies written as its
terms with explicit annotations, type-checked at compile time and translated by
\(\sem{-}\). An underfunded comprehension does not complete, and the deployer
is notified through the budget trace.
\item Stage~1 aggregation: \texttt{by} a theory instance descending from
\texttt{CommutativeMonoid}, checked by origin and defaulting to
\texttt{ParMonoid}, together with its two diagnostics.
\item The cost-accounted surface with both gates, per-binding outer charge, and
the \(n\)-ary outer gate for single-group comprehensions.
\item \fe{} on the left of \texttt{;}.
\item Native evaluation in the reducer, the new process node gated on protocol
version, and the demand rule in the funding analysis.
\item The laws, the bag morphism, confluence, the fragment obligation and the
aggregator obligations recorded, and discharged for the prelude monads and for
\texttt{ParMonoid}.
\end{enumerate}

\noindent
Next: the second terminating fragment, Abramsky's proof expressions.
Deferred: checking hand-written \flang{} bodies directly against
\(\varphi_\tau\) with a type system generated by the Hypercube construction;
sort parameters and user-declared monads (Step~2), with the
parametric \rhoc{} calculus of \S\ref{sec:rho} as their acceptance example;
syntax monads declared by substitution; composition through distributive laws;
aggregation Stages~2--4; graded \texttt{where} in \fe{}; the lax modality.

%=============================================================================
\section{Open questions}
\label{sec:questions}

\begin{question}[The logic of the fragment]
\label{q:logic}
Part (iii) of Obligation~\ref{obl:fragment} needs least fixed points in the
generated logics, and for System F quantification over formulae. Should the
generated logics acquire both, so that the fragment can be stated in them and,
eventually, hand-written bodies checked against it; or should \(\varphi_\tau\)
target a fixed family of formulae and stay first order?
\end{question}

\begin{question}[Element sorts in System F]
A body's pattern variables have the monad's element sort, and that sort must be
the translation of a System F type. Should a theory's sorts be imported into the
System F theory as base types, or represented by Church encodings, which System
F supports but which change what the body sees?
\end{question}

\begin{question}[Dependent groups]
\label{q:dependent}
For a \fe{} with dependent groups, should the outer gates be taken group by
group as the deficit of \S\ref{sec:refsem} unfolds, accepting that a later group
may find its lane empty after earlier groups have paid, or should a declared
upper bound be taken at entry, accepting overpayment when the bound is loose?
\end{question}

\begin{question}[Per-element funding]
\label{q:hash}
May the inner signature mention variables bound by the pattern,
\verb|{% P %}[#x]|? It lets data pay for its own processing, at the price of
data-dependent lanes in the deployer's estimate.
\end{question}

\begin{question}[Semilattice in the tower]
Stage~2 needs a \texttt{Semilattice} theory, and Stage~4 a \texttt{Semimodule}
over a rig. Should both join the universal-algebra module now, so that the
origin checks of later stages have something to descend from?
\end{question}

\begin{question}[Arity drift]
Whatever is decided for arity-changing replacements of \texttt{PPar} must be
decided for monad interpretations and for aggregator operations. Which
resolution is preferred?
\end{question}

\begin{question}[Syntax monads]
Should every theory with a variable former get its syntax monad, declared
\texttt{by substitution at} that former, with laws discharged once for the
toolchain, and with \(\mathsf{toBag}\) always taken on the node's canonical
form?
\end{question}

\begin{question}[Generated distributive laws]
For which group actions on a theory can the toolchain generate the
distributive law, as it can for the deferred action of oriented processes, and
how should a law that cannot be generated be supplied and its four obligations
stated?
\end{question}

\begin{question}[Default monad]
Should a sort with exactly one declared monad allow the qualifier to be omitted,
as proposed, or should the qualifier always be required for readability?
\end{question}

\begin{thebibliography}{99}\small\raggedright

\bibitem{Abramsky94}
S.~Abramsky.
\newblock Proofs as processes.
\newblock \emph{Theoretical Computer Science}, 135(1):5--9, 1994.

\bibitem{BentonBiermanDePaiva98}
P.~N.~Benton, G.~M.~Bierman, and V.~C.~V.~de~Paiva.
\newblock Computational types from a logical perspective.
\newblock \emph{Journal of Functional Programming}, 8(2):177--193, 1998.

\bibitem{CairesPfenning10}
L.~Caires and F.~Pfenning.
\newblock Session types as intuitionistic linear propositions.
\newblock In \emph{CONCUR 2010}, LNCS 6269, pages 222--236, 2010.

\bibitem{DeanGhemawat04}
J.~Dean and S.~Ghemawat.
\newblock MapReduce: simplified data processing on large clusters.
\newblock In \emph{Proceedings of the 6th USENIX Symposium on Operating Systems
Design and Implementation (OSDI)}, pages 137--150, 2004.

\bibitem{FairtloughMendler97}
M.~Fairtlough and M.~Mendler.
\newblock Propositional lax logic.
\newblock \emph{Information and Computation}, 137(1):1--33, 1997.

\bibitem{AltenkirchReus99}
T.~Altenkirch and B.~Reus.
\newblock Monadic presentations of lambda terms using generalized inductive
types.
\newblock In \emph{Computer Science Logic (CSL)}, LNCS 1683, pages 453--468,
1999.

\bibitem{Beck69}
J.~Beck.
\newblock Distributive laws.
\newblock In \emph{Seminar on Triples and Categorical Homology Theory}, Lecture
Notes in Mathematics 80, pages 119--140. Springer, 1969.

\bibitem{BirdPaterson99}
R.~S.~Bird and R.~Paterson.
\newblock De Bruijn notation as a nested datatype.
\newblock \emph{Journal of Functional Programming}, 9(1):77--91, 1999.

\bibitem{FiorePlotkinTuri99}
M.~Fiore, G.~Plotkin, and D.~Turi.
\newblock Abstract syntax and variable binding.
\newblock In \emph{Proceedings of the 14th IEEE Symposium on Logic in Computer
Science (LICS)}, pages 193--202, 1999.

\bibitem{GirardLafontTaylor89}
J.-Y.~Girard, Y.~Lafont, and P.~Taylor.
\newblock \emph{Proofs and Types}.
\newblock Cambridge Tracts in Theoretical Computer Science 7. Cambridge
University Press, 1989.

\bibitem{Kock70}
A.~Kock.
\newblock Monads on symmetric monoidal closed categories.
\newblock \emph{Archiv der Mathematik}, 21:1--10, 1970.

\bibitem{LINQ}
E.~Meijer, B.~Beckman, and G.~Bierman.
\newblock LINQ: reconciling object, relations and XML in the .NET framework.
\newblock In \emph{Proceedings of the ACM SIGMOD International Conference on
Management of Data}, page 706, 2006.

\bibitem{Moggi91}
E.~Moggi.
\newblock Notions of computation and monads.
\newblock \emph{Information and Computation}, 93(1):55--92, 1991.

\bibitem{Reynolds74}
J.~C.~Reynolds.
\newblock Towards a theory of type structure.
\newblock In \emph{Programming Symposium}, LNCS 19, pages 408--425, 1974.

\bibitem{Wadler87}
P.~Wadler.
\newblock List comprehensions.
\newblock In S.~L.~Peyton~Jones, \emph{The Implementation of Functional
Programming Languages}, chapter~7. Prentice Hall, 1987.

\bibitem{Wadler92}
P.~Wadler.
\newblock Comprehending monads.
\newblock \emph{Mathematical Structures in Computer Science}, 2(4):461--493,
1992.

\bibitem{Wadler12}
P.~Wadler.
\newblock Propositions as sessions.
\newblock In \emph{Proceedings of the 17th ACM SIGPLAN International Conference
on Functional Programming (ICFP)}, pages 273--286, 2012.

\bibitem{Wells99}
J.~B.~Wells.
\newblock Typability and type checking in System F are equivalent and
undecidable.
\newblock \emph{Annals of Pure and Applied Logic}, 98(1--3):111--156, 1999.

\bigskip
\noindent\textbf{F1R3FLY.io working notes.} Each lives in the F1R3FLY.io
\texttt{publications} repository, in the directory given.

\bibitem{CostRho}
L.~G.~Meredith.
\newblock Cost-accounted rho calculus: a spectral decomposition of phlogiston.
\newblock F1R3FLY.io working note, 2026. \texttt{cost-accounting/}.

\bibitem{CostMonad}
L.~G.~Meredith.
\newblock Continued interactive GSLTs and the cost endofunctor.
\newblock F1R3FLY.io working note, 2026. \texttt{cost-accounting-as-monad/}.

\bibitem{F1R3Dev}
L.~G.~Meredith.
\newblock \flang{} for the working software developer.
\newblock F1R3FLY.io working note, 2026.
\texttt{MeTTaIL4WorkingDev/f1r3lang4dev/}.

\bibitem{GradedWhere}
L.~G.~Meredith.
\newblock Graded where-clauses: non-crisp truth values as the syntax of
non-determinism resolution in the rho calculus.
\newblock F1R3FLY.io working note, 2026. \texttt{fuzzyware/}.

\bibitem{GSLTintro}
L.~G.~Meredith.
\newblock Graph-structured lambda theories: a reference for the working
developer.
\newblock F1R3FLY.io working note, 2026. \texttt{GSLT-intro/}.

\bibitem{History}
L.~G.~Meredith.
\newblock History as an endofunctor: reversibility, covering spaces, and the
memory price of energy conservation in cost-accounted GSLTs.
\newblock F1R3FLY.io working note, 2026. \texttt{history-functor/}.

\bibitem{TwoErasures}
L.~G.~Meredith.
\newblock Two erasures and two laxities: internalising the cost and history
monads over continued interactive GSLTs.
\newblock F1R3FLY.io working note, 2026. \texttt{history-and-cost/}.

\bibitem{Namespaces}
L.~G.~Meredith.
\newblock Namespaces without a name server: freshness by communication, and a
fully abstract encoding of the \(\pi\)-calculus in the rho calculus.
\newblock F1R3FLY.io working note, 2026. \texttt{pi2rhobugfix/}.

\bibitem{Oriented}
L.~G.~Meredith.
\newblock Oriented processes: group-valued data, a chaos constant, and an
equivariant COMM for the reflective higher-order calculus.
\newblock F1R3FLY.io working note, draft~3, 2026. \texttt{process-rotations/}.

\bibitem{PF}
L.~G.~Meredith.
\newblock Plausible fiction on \flang{}: a requirements and design document for
a typed marketplace of programmatic intentions.
\newblock F1R3FLY.io working note, 2026. \texttt{plausible-fiction/}.

\bibitem{WSH}
L.~G.~Meredith.
\newblock Where substitution happens: semantic substitution, deferred unquote
and explicit substitution in a paired-message rho calculus.
\newblock F1R3FLY.io working note, 2026. \texttt{OnSemanticSubstitution/}.

\bibitem{Hypercube}
M.~Stay, L.~G.~Meredith, and J.~B.~Wells.
\newblock Generating hypercubes of type systems.
\newblock F1R3FLY.io working note, 2026. \texttt{drafts/Hypercube/}.

\bibitem{WeightedGSLT}
L.~G.~Meredith.
\newblock Weighted graph-structured lambda theories: stochastic and quantum
execution for \flang{}.
\newblock F1R3FLY.io working note, second version, 2026. \texttt{rho-life/}.

\bibitem{MR05}
L.~G.~Meredith and M.~Radestock.
\newblock A reflective higher-order calculus.
\newblock \emph{Electronic Notes in Theoretical Computer Science},
141(5):49--67, 2005.

\bigskip
\noindent\textbf{Software.}

\bibitem{modulesyntax}
F1R3FLY.io.
\newblock \texttt{f1r3node-rust}, branch \texttt{feature/module-syntax}: the
\texttt{mettail-elab} reference elaborator for the \texttt{Module} and
\texttt{Theory} surface.
\newblock \url{https://github.com/F1R3FLY-io/f1r3node-rust}.

\bibitem{f1r3nodeOTE}
F1R3FLY.io.
\newblock \texttt{f1r3node-rust}, branch \texttt{feature/f1r3lang-mettail-only}:
the frontend admission contract (\texttt{docs/design/f1r3lang-frontend.md}),
the evaluation composition (\texttt{docs/node/f1r3lang-evaluation.md}), the
cost-monad correspondence record and graded adequacy development
(\texttt{docs/theory/}), the synchronous-send normaliser, and the funding
analysis (\texttt{rholang/src/rust/interpreter/accounting/delta\_sigma.rs}).
\newblock \url{https://github.com/F1R3FLY-io/f1r3node-rust}.

\end{thebibliography}

\end{document}
